\documentclass[12pt,a4paper]{article}
\usepackage[utf8]{inputenc}
\usepackage[T1]{fontenc}
\usepackage{amsmath,amssymb,amsfonts}
\usepackage{graphicx}
\usepackage{booktabs}
\usepackage{xcolor}
\usepackage[
    colorlinks=true,
    linkcolor=blue!70!black,
    citecolor=green!50!black,
    urlcolor=blue!70!black,
    bookmarks=true,
    bookmarksopen=true,
    pdfauthor={Razvan-Constantin Anghelina},
    pdftitle={Geometric Foundations of Universal Constants: A 600-Cell Framework},
    pdfsubject={Theoretical Physics},
    pdfkeywords={600-cell, fine structure constant, gauge unification, golden ratio, Planck scale}
]{hyperref}
\usepackage{geometry}
\usepackage{float}
\usepackage{caption}
\usepackage{subcaption}
\usepackage{multirow}

\geometry{margin=2.5cm}

\title{\textbf{Geometric Foundations of Universal Constants:\\
A 600-Cell Framework for Coupling Constants, Mass Hierarchies, and Spacetime Geometry}}

\author{
Razvan-Constantin Anghelina\\
\texttt{contact@webdesignmedia.ro}\\[1em]
\textit{Version 3.6 -- February 2026}
% v3.4: Edge split 1+3+8, Gram=24-cell, A_eff=A^2/3 (EXP-143,144)
}

\date{}

\begin{document}

\maketitle

\begin{abstract}
We present a framework in which the fundamental constants of particle physics emerge from the geometric and spectral properties of the 600-cell polytope. This four-dimensional regular polytope, identified as the universal energy minimizer in $\mathbb{R}^4$ by the Cohn-Kumar theorem, provides a discrete structure whose spectral properties encode the coupling constants with high precision: the fine structure constant $\alpha$ (0.0001\%), the strong coupling $\alpha_s$ (0.03\%), the Weinberg angle $\sin^2\theta_W$ (0.19\%), and the Higgs mass (0.09\%). The 24 non-fermionic vertices form the $D_4$ root system, from which the Standard Model gauge group $SU(3) \times SU(2) \times U(1)$ emerges as a maximal-rank subalgebra. The remaining 96 vertices of the snub 24-cell encode fermions ($96 = 16 \times 3 \times 2$), with a geometric 4-to-3 selection mechanism via the temporal axis. \textbf{The lepton mass hierarchy} $m_\mu/m_e \sim \phi^{11}$, $m_\tau/m_\mu \sim \phi^{6}$ \textbf{is derived from holonomy on the Hopf fibration}, where the exponents 5 (graph diameter) and 6 (decagons per vertex) are topological invariants of the 600-cell. Electroweak phase corrections reduce the lepton mass errors to 0.25--0.53\%. A unified mass formula $m_f = m_e \cdot \phi^{a \cdot d_{\text{eff}} + b \cdot k_{\text{eff}}}$ with sector-dependent corrections predicts seven of nine fermion masses within 2\%. We identify Discrete Scale Invariance (DSI) as the physical mechanism: the 600-cell lattice generates a potential with minima at $m_e \cdot \phi^n$, and the topological quantum numbers $(a,b)$ count diameter traversals and decagonal loops on the graph. An exploratory gravitational model (Surface Tension Gravity) is presented in the appendix: it reproduces GR weak-field tests at 1PN order but predicts a black hole shadow half the GR size ($\sim 7\sigma$ tension with EHT), and should be regarded as a toy model. The topological quantum numbers $(a,b)$ are uniquely determined by complexity minimization on the graph: $(a,b) = \arg\min|a|+|b|$ subject to $5a + 6b = n$, where 5 and 6 are the intersection numbers of the 600-cell. The $E_8$ root system is constructed explicitly from the 600-cell via the icosian lattice: two copies $S$ and $T = \phi' \cdot S$ yield all 240 roots with exact inner products. Two algebraic constraints---$\mathbb{Z}[\phi]$ norm restriction and a three-line rule in the $(a,b)$ lattice---plus norm positivity on the $b=1$ line correctly classify all 31 mass levels from $n=0$ to $n=30$; the sole remaining pattern is the three-generation limit ($b \leq 2$). The CKM exponents connect to $H_4$ Coxeter exponents and admit a cleaner formulation via $H_4 \times H_4'$ cross-pairings from the $E_8$ branching. The Down quark prediction (5.15\,MeV) lies within $1\sigma$ of the PDG value ($4.67^{+0.48}_{-0.17}$\,MeV). While $D_4$ triality is preserved by the 600-cell geometry alone, the $E_8$ embedding breaks it: the three $D_4$ legs extend into $E_8$ with distinct tail lengths $\{0, 1, 3\}$ and Kac labels $\{3, 4, 5\}$, uniquely assigning each leg to $SU(3)_C$, $SU(2)_L$, or $U(1)_Y$ via the standard GUT breaking chain. The edge structure automatically forbids direct gauge--gauge couplings (A-A $=$ A-B $=$ B-B $= 0$), leaving exactly three interaction vertex types (ACC, BCC, CCC) in the ratio $1:2:2$; the massless graph propagator changes sign at distance $d=3$, exhibiting confinement-like behavior. A detailed analysis of the local neighborhood reveals a natural $1 + 3 + 8$ edge splitting that matches $\dim(U(1) \times SU(2) \times SU(3)) = 12$: each fermion has a geometrically unique ``special'' CC neighbor (degree~4, shared A-gauge vertex, zero shared B), and the effective gauge adjacency through fermion mediators reproduces the 24-cell exactly ($A_{\text{eff}} = A^2|_{\text{gauge}}/3$). Among 1200 plaquettes, only $SU(3)$ has pure plaquettes (288)---zero for $U(1)$ and $SU(2)$---reproducing the non-Abelian self-interaction structure of QCD. A vertex-transitivity theorem proves that spectral projections are sector-blind, establishing that the Lagrangian \emph{form} (from topology) and coupling \emph{values} (from the spectrum) are two independent geometric properties of the 600-cell. We explicitly identify remaining open problems including the three-generation count, the CKM derivation mechanism, and the need for a complete Lagrangian formulation.
\end{abstract}

\tableofcontents
\newpage

%==============================================================
\section{Introduction}
%==============================================================

The search for a fundamental structure underlying the constants of nature has been a central pursuit in theoretical physics. The Standard Model of particle physics, while remarkably successful, contains approximately 19 free parameters whose values must be determined experimentally. The origin of these numbers---why the fine structure constant equals $1/137.036$ rather than some other value---remains one of the deepest mysteries in physics.

We propose that spacetime at the Planck scale possesses the discrete structure of the 600-cell (hexacosichoron), a regular four-dimensional polytope with exceptional geometric properties. This approach is motivated by the Cohn-Kumar theorem \cite{cohn2007,cohn2011}, which establishes the 600-cell as the unique universal energy minimizer in $\mathbb{R}^4$ for a broad class of potentials. The geometric framework builds upon the 600-cell GUT developed by O'Neill \cite{oneill2025gut,oneill2023binary}.

The key insight is that the fundamental constants are not arbitrary but emerge as ``spectral fingerprints'' of this optimal geometric structure. The golden ratio $\phi = (1+\sqrt{5})/2$, which permeates the 600-cell geometry, appears naturally in the coupling constants and mixing angles of particle physics.

This paper is organized as follows: Section~2 describes the 600-cell architecture and its physical interpretation. Section~3 derives the fundamental coupling constants from the Laplacian spectrum. Section~4 addresses the Higgs sector and mass generation. Section~5 maps fermions to the vertex structure, derives the Standard Model gauge group from the $D_4$ root system, and explains the 4-to-3 generation selection mechanism. Section~6 presents the CKM and PMNS mixing matrices. Section~7 derives the fermion mass hierarchy from holonomy on the Hopf fibration, culminating in a unified mass formula for all nine fermions. Section~8 identifies Discrete Scale Invariance as the physical mechanism underlying the mass formula. Section~9 constructs the $E_8$ root system from the icosian lattice and discusses gauge unification. Section~10 validates the Lagrangian framework. Section~11 proposes experimental tests, and Sections~12--13 summarize results and conclude with an explicit discussion of limitations and open problems. Appendix~A presents an exploratory gravitational model (Surface Tension Gravity) that illustrates how the 600-cell geometry may constrain spacetime structure, while noting its tension with strong-field observations.

%==============================================================
\section{The 600-Cell Architecture}
%==============================================================

\subsection{Fundamental Properties}

The 600-cell is a regular convex 4-polytope with the following characteristics:

\begin{table}[H]
\centering
\begin{tabular}{lrl}
\toprule
\textbf{Property} & \textbf{Value} & \textbf{Physical Interpretation} \\
\midrule
Vertices & 120 & Lattice nodes / degrees of freedom \\
Edges & 720 & Propagation paths \\
Faces & 1200 & Triangular interfaces \\
Cells & 600 & Tetrahedral units \\
Symmetry group & $H_4$ & Order 14,400 \\
Vertex degree & 12 & Neighbors per node \\
Graph diameter & 5 & Maximum geodesic distance \\
Euler characteristic & 0 & Topological invariant \\
\bottomrule
\end{tabular}
\caption{Properties of the 600-cell polytope.}
\end{table}

\subsection{Universal Optimality}

\textbf{Theorem (Cohn-Kumar \cite{cohn2007}):} The 600-cell is the unique configuration in $\mathbb{R}^4$ that acts as a universal energy minimizer for completely monotonic potentials.

\textbf{Physical consequence:} By the principle of least action, spacetime at the Planck scale self-organizes into the 600-cell lattice structure, analogous to:
\begin{itemize}
    \item 2D: Hexagonal lattice (honeycomb)
    \item 3D: FCC/HCP crystal structures
    \item 4D: 600-cell
\end{itemize}

\subsection{Vertex Decomposition}

The 120 vertices decompose into three geometrically distinct sets:

\begin{equation}
120 = 8 + 16 + 96
\end{equation}

\begin{table}[H]
\centering
\begin{tabular}{ccll}
\toprule
\textbf{Set} & \textbf{Count} & \textbf{Structure} & \textbf{Physical Role} \\
\midrule
Type A & 8 & 16-cell (cross-polytope) & \multirow{2}{*}{$D_4$ root system $\to$ gauge sector (Section~\ref{sec:gauge})} \\
Type B & 16 & 8-cell (tesseract) & \\
Type C & 96 & Snub 24-cell & Fermions ($16 \times 3 \times 2$) \\
\bottomrule
\end{tabular}
\caption{Vertex decomposition and physical interpretation.}
\end{table}

The 96 golden-ratio vertices encode:
\begin{equation}
96 = 16 \times 3 \times 2 = \text{(fermions per generation)} \times \text{(generations)} \times \text{(chiralities)}
\end{equation}

\subsection{Spectral Structure and Association Scheme}
\label{sec:spectral}

The 600-cell graph is \textbf{distance-transitive}: vertices at the same graph distance are related by symmetry. This creates an \textbf{association scheme} with 9 angular distance classes $d \in \{0, 1, \ldots, 8\}$ (corresponding to $S^3$ arc distances $\theta = 0, \pi/5, \ldots, \pi$), and the adjacency matrix has exactly 9 distinct eigenvalues---all of the form $a + b\phi$ with $a,b \in \mathbb{Z}$:

\begin{table}[H]
\centering
\begin{tabular}{cccl}
\toprule
\textbf{Index} $i$ & \textbf{Eigenvalue} $\theta_i$ & \textbf{Multiplicity} $m_i$ & \textbf{Structure} \\
\midrule
0 & 12             & $1 = 1^2$   & Trivial representation \\
1 & $6\phi$        & $4 = 2^2$   & \multirow{2}{*}{Galois pair in $\mathbb{Q}(\sqrt{5})$} \\
8 & $6 - 6\phi$    & $4 = 2^2$   & \\
2 & $4\phi$        & $9 = 3^2$   & \multirow{2}{*}{Galois pair} \\
6 & $4 - 4\phi$    & $9 = 3^2$   & \\
3 & 3              & $16 = 4^2$  & \multirow{2}{*}{Rational pair} \\
7 & $-3$           & $16 = 4^2$  & \\
4 & 0              & $25 = 5^2$  & Kernel \\
5 & $-2$           & $36 = 6^2$  & Largest eigenspace \\
\bottomrule
\end{tabular}
\caption{Adjacency spectrum of the 600-cell. All multiplicities are perfect squares $n^2$ for $n = 1, \ldots, 6$, reflecting $H_4$ irreducible representations.}
\label{tab:spectrum}
\end{table}

The characteristic polynomial over $\mathbb{Q}$ factors as:
\begin{equation}
P(x) = (x-12)(x^2-6x-36)^4(x^2-4x-16)^9(x-3)^{16}\,x^{25}(x+2)^{36}(x+3)^{16}
\end{equation}

Two key structural features:
\begin{enumerate}
    \item \textbf{Intersection numbers:} Each edge of the 600-cell has exactly $a_1 = 5$ common neighbors and $b_1 = 6$ next-shell neighbors. These are topological invariants---identical for all 720 edges.
    \item \textbf{Phi-sector:} The four eigenspaces with $\phi$-dependent eigenvalues ($i = 1, 2, 6, 8$) have total dimension $4 + 9 + 9 + 4 = 26$. This is the same 26 that appears in $\sin^2\theta_W = 6/26$ (Section~3.3).
\end{enumerate}

Three exact \textbf{Laplacian eigenvalue ratios} connect the spectrum to physics:
\begin{align}
\lambda_4 / \lambda_1 &= 2\phi^2 &\quad &\text{(yields } \alpha \text{, Section~3.1)} \\
\lambda_6 / \lambda_2 &= \phi^2 &\quad &\text{(connects to QED vertex corrections)} \\
\lambda_7 / \lambda_3 &= 5/3 = F/E &\quad &\text{(faces/edges = plaquettes per link, EXP-132)}
\end{align}

where $\lambda_i = 12 - \theta_i$ are Laplacian eigenvalues.

%==============================================================
\section{Derivation of Fundamental Constants}
%==============================================================

\subsection{The Fine Structure Constant}

The fine structure constant $\alpha$ emerges from a consistency equation relating the Laplacian spectrum to geometric phases:

\begin{equation}
\boxed{2\pi\alpha^2 - 20\phi^4\alpha + 1 = 0}
\end{equation}

\textbf{Origin of terms:}
\begin{itemize}
    \item $20\phi^4 = 5 \times (\lambda_4/\lambda_1)^2$ from the 600-cell Laplacian spectrum
    \item $\lambda_4/\lambda_1 = 2\phi^2$ (exact eigenvalue ratio)
    \item $2\pi$ = geometric phase accumulated on a Hopf fiber (decagon loop)
\end{itemize}

\textbf{Solution:}
\begin{equation}
\alpha = \frac{20\phi^4 - \sqrt{(20\phi^4)^2 - 8\pi}}{4\pi} = \frac{1}{137.036...}
\end{equation}

\begin{table}[H]
\centering
\begin{tabular}{lcc}
\toprule
& \textbf{Calculated} & \textbf{Experimental \cite{codata2022}} \\
\midrule
$\alpha^{-1}$ & 137.036 & $137.035999084(21)$ \\
\textbf{Error} & \multicolumn{2}{c}{\textbf{0.0001\%}} \\
\bottomrule
\end{tabular}
\caption{Fine structure constant comparison.}
\end{table}

\subsection{The Strong Coupling Constant}

The strong coupling constant at the $Z$ mass scale:

\begin{equation}
\alpha_s(M_Z) = \frac{1}{2\phi^3} = 0.1180
\end{equation}

\textbf{Exact algebraic relation:}
\begin{equation}
\boxed{\frac{\alpha_s}{\alpha} = 10\phi}
\end{equation}

This ratio has a spectral derivation. From Table~\ref{tab:spectrum}, the eigenvalue ratio $\theta_1/\theta_3 = 6\phi/3 = 2\phi$, and the intersection number $a_1 = 5$ (Section~\ref{sec:spectral}), so:
\begin{equation}
\frac{\alpha_s}{\alpha} = a_1 \cdot \frac{\theta_1}{\theta_3} = 5 \times 2\phi = 10\phi
\end{equation}

\textbf{Classification:} DERIVED---the ratio $\alpha_s/\alpha$ is determined by the first intersection number ($a_1 = 5$) and an exact eigenvalue ratio ($\theta_1/\theta_3 = 2\phi$), both topological invariants of the 600-cell.

\begin{table}[H]
\centering
\begin{tabular}{lcc}
\toprule
& \textbf{Calculated} & \textbf{Experimental} \\
\midrule
$\alpha_s(M_Z)$ & 0.1180 & $0.1179 \pm 0.0009$ \\
\textbf{Error} & \multicolumn{2}{c}{\textbf{0.03\%}} \\
\bottomrule
\end{tabular}
\caption{Strong coupling constant comparison.}
\end{table}

\subsection{The Weinberg Angle}

The electroweak mixing angle derives from geometric structure counts:

\begin{equation}
\sin^2\theta_W = \frac{6}{26} = \frac{6}{6+20} = 0.2308
\end{equation}

\textbf{Geometric interpretation:}
\begin{itemize}
    \item 6 = decagonal directions per vertex (U(1) sector)
    \item 20 = tetrahedral cells per vertex (SU(2) sector)
    \item 26 = total structures
\end{itemize}

\textbf{Second derivation path (spectral):} The same ratio emerges from the association scheme (Section~\ref{sec:spectral}):
\begin{equation}
\sin^2\theta_W = \frac{b_1}{\dim(\phi\text{-sector})} = \frac{6}{26}
\end{equation}
where $b_1 = 6$ is the second intersection number and $\dim(\phi\text{-sector}) = 4 + 9 + 9 + 4 = 26$ is the total dimension of the four $\phi$-dependent eigenspaces. Two independent geometric arguments yielding the same fraction strengthens the derivation.

\begin{table}[H]
\centering
\begin{tabular}{lcc}
\toprule
& \textbf{Calculated} & \textbf{Experimental} \\
\midrule
$\sin^2\theta_W$ & 0.2308 & 0.2312 \\
\textbf{Error} & \multicolumn{2}{c}{\textbf{0.19\%}} \\
\bottomrule
\end{tabular}
\caption{Weinberg angle comparison.}
\end{table}

%==============================================================
\section{The Higgs Sector}
%==============================================================

\subsection{Mass from Dihedral Angles}

The Higgs mass represents the phase transition energy from tetrahedral (600-cell) to octahedral (24-cell) geometry:

\textbf{Bare mass:}
\begin{equation}
m_H^{(0)} = m_W \cdot \frac{\theta_{\text{octahedron}}}{\theta_{\text{tetrahedron}}} = m_W \cdot \frac{\arccos(-1/3)}{\arccos(1/3)} = 124.76 \text{ GeV}
\end{equation}

\textbf{Topological correction:}
The 8 vertices of the inscribed 16-cell contribute a correction of $-8\alpha$:

\begin{equation}
\boxed{m_H = m_W \cdot (\phi - 8\alpha) = 125.36 \text{ GeV}}
\end{equation}

\begin{table}[H]
\centering
\begin{tabular}{lcc}
\toprule
& \textbf{Calculated} & \textbf{Experimental (LHC)} \\
\midrule
$m_H$ & 125.36 GeV & 125.25 GeV \\
\textbf{Error} & \multicolumn{2}{c}{\textbf{0.09\%}} \\
\bottomrule
\end{tabular}
\caption{Higgs mass comparison.}
\end{table}

\subsection{The Higgs-STG Portal}

Following the Surface Tension Gravity framework \cite{oneill2025stg,oneill2025gps,oneill2025scaffold} (see Appendix~\ref{app:stg}), the Lagrangian includes a portal coupling between the Higgs field and the gravitational compression factor:

\begin{equation}
\mathcal{L}_{\text{portal}} = \lambda_{CH}^2 C^2 (H^\dagger H)
\end{equation}

where $C(r) = 1 + \alpha_G/r$ is the gravitational compression factor.

%==============================================================
\section{Fermion Mapping}
%==============================================================

\subsection{Three Generations from Geometry}

The 96 vertices of the snub 24-cell (golden-ratio vertices) accommodate exactly three generations of fermions:

\begin{equation}
96 = 16 \times 3 \times 2
\end{equation}

\begin{itemize}
    \item 16 = Weyl fermions per generation (quarks + leptons with color)
    \item 3 = generations
    \item 2 = chiralities (L/R)
\end{itemize}

\subsection{The 4-to-3 Generation Selection Mechanism}

A deeper analysis (EXP-090) reveals that the 96 vertices naturally decompose into \textbf{four groups of 24}, based on which coordinate vanishes:
\begin{equation}
96 = 4 \times 24
\end{equation}

\textbf{Key observation:} Each axial direction $k$ in $\mathbb{R}^4$ connects only to the groups where $x_k \neq 0$. Therefore:

\begin{itemize}
    \item Choosing a \textbf{privileged axis} (the time direction, e.g., $x_3$) \textbf{selects 3 groups}
    \item The fourth group (where $x_3 = 0$) becomes \textbf{disconnected} from the temporal axis
    \item This group is identified with a \textbf{sterile sector} (right-handed neutrinos, dark matter candidates)
\end{itemize}

\textbf{Physical interpretation:} The breaking of the $\mathbb{R}^4$ symmetry by the choice of time direction naturally reduces 4 geometric groups to 3 observable generations plus 1 sterile sector.

\subsection{The Gauge Sector: $D_4$ Root System from the 24-Cell}
\label{sec:gauge}

The 24 non-fermionic vertices (Type~A + Type~B) form a \textbf{24-cell}, which is simultaneously a regular 4-polytope and a root system. This root system is $D_4$.

\textbf{Verification.} The 8 Type-A vertices $(\pm 1, 0, 0, 0)$, etc., together with the 16 Type-B vertices $(\pm\tfrac{1}{2})^4$, satisfy the Cartan integer condition $2\langle \alpha, \beta\rangle / \langle \alpha, \alpha\rangle \in \mathbb{Z}$ for all root pairs, confirming they form the $D_4$ root system (rank 4, 24 roots).

\textbf{Standard Model gauge group.} The maximal-rank regular subalgebra decomposition of $D_4$ gives:
\begin{equation}
\boxed{D_4 \;\to\; A_2 \oplus A_1 \oplus U(1) \;=\; SU(3) \times SU(2) \times U(1)}
\end{equation}

This is a standard result in Lie theory. The dimension and rank match exactly:
\begin{align}
\dim(SU(3) \times SU(2) \times U(1)) &= 8 + 3 + 1 = 12 = \text{vertex degree of 600-cell} \\
\mathrm{rank}(SU(3) \times SU(2) \times U(1)) &= 2 + 1 + 1 = 4 = \dim(\mathbb{R}^4)
\end{align}

The remaining 96 vertices (Type~C, the snub 24-cell) naturally separate as the \textbf{matter sector}, with $96 = 16 \times 3 \times 2$ encoding 16 fermion states per generation, 3 generations, and 2 chiralities.

\textbf{$D_4$ triality.} The Dynkin diagram of $D_4$ has an $S_3$ outer automorphism group (triality), connecting three inequivalent 8-dimensional representations: $\mathbf{8}_v$, $\mathbf{8}_s$, $\mathbf{8}_c$. Left quaternion multiplication by $\omega = \tfrac{1}{2}(1,1,1,1)$ is a full 600-cell automorphism that cyclically permutes $\mathbf{8}_v \to \mathbf{8}_s \to \mathbf{8}_c \to \mathbf{8}_v$, while preserving the 96 Type-C (matter) vertices. Consequently, the apparent $8 + 16$ split (Type~A versus Type~B) is a \emph{coordinate artifact}: the triality symmetry is \textbf{fully preserved} by the 600-cell geometry (EXP-126).

\textbf{$A_5$ irrep selection.} Although the 600-cell alone cannot break triality, the \emph{vertex figure} (icosahedron, symmetry $A_5$) provides a partial selection. The 12 edges at each vertex decompose under $A_5$ as $12 = 1 + 3 + 3' + 5$. The \emph{unique} regrouping compatible with SM dimensions is $8 = 3' + 5$, $3 = 3$, $1 = 1$, corresponding to $SU(3) \times SU(2) \times U(1)$. The competing decomposition $SU(2)^4$ would require $3 + 3 + 3 + 3$, which is incompatible with the $A_5$ structure.

\textbf{$E_8$ triality breaking (EXP-130, 133).} The $E_8$ embedding (Section~\ref{sec:e8}) resolves the triality question definitively. The triality automorphism $\omega$ maps all 240 $E_8$ roots to $E_8$ roots and preserves both 600-cell copies $S$ and $T$ separately---so at the \emph{root system} level, triality is an $E_8$ symmetry. However, at the \emph{Dynkin diagram} level, triality is \textbf{broken}. The $D_4$ sub-diagram $\{2, 3, [4], 5\}$ sits inside the $E_8$ Dynkin diagram (Bourbaki numbering):
\begin{equation}
1 - 3 - \boxed{4} - 5 - 6 - 7 - 8 \qquad \text{with} \quad 2 - \boxed{4}
\end{equation}
where node $[4]$ is the $D_4$ branching node. The three legs extend into $E_8$ with \textbf{unequal tail lengths}, and the Kac labels (marks of the highest root) are all \textbf{distinct}:

\begin{table}[H]
\centering
\begin{tabular}{ccccl}
\toprule
\textbf{D4 Leg} & \textbf{Tail length} & \textbf{Kac label} & \textbf{Tail algebra} & \textbf{SM assignment} \\
\midrule
Node 5 $\to$ 6,7,8 & 3 & $m = 5$ & $A_4 \supset A_2$ & $SU(3)_{\text{color}}$ \\
Node 3 $\to$ 1       & 1 & $m = 4$ & $A_1$              & $SU(2)_{\text{weak}}$ \\
Node 2 (dead end)     & 0 & $m = 3$ & ---                & $U(1)_Y$ (via $SU(2)_R$) \\
\bottomrule
\end{tabular}
\caption{$D_4$ leg assignment from $E_8$ Dynkin diagram topology. Tail length: number of $E_8$ nodes beyond $D_4$ on each leg.}
\end{table}

The assignment is \textbf{unique and monotonic}: longer tail $\to$ larger gauge group $\to$ more generators. Since $E_8$ has no nontrivial outer automorphisms, the triality can only act as an inner automorphism (Weyl group element), which assigns distinct roles to the three legs via the Weyl chamber structure. The standard GUT breaking chain is fully recovered:
\begin{equation}
E_8 \;\to\; E_6 \times SU(3)_{\text{fam}} \;\to\; SO(10) \;\to\; SU(5) \;\to\; SU(3)_C \times SU(2)_L \times U(1)_Y
\end{equation}

The $D_4$ simple roots within $E_8$ come from \emph{mixed} 600-cells: 2 from $S_{96}$ and 2 from $T_{96}$ (snub 24-cell vertices of each icosian copy), confirming that the two 600-cells play asymmetric roles in the $E_8$ root structure.

\textbf{$SU(2)_R \to U(1)_Y$ breaking (EXP-135).} The dead-end leg (node 2) generates $SU(2)_R$, which reduces to $U(1)_Y$ through three independent structural mechanisms:
\begin{enumerate}
    \item \textbf{Topological isolation:} The dead end has no non-Abelian extension in $E_8$ to protect $SU(2)_R$ from breaking.
    \item \textbf{Maximal commutant:} Removing node 2 from $E_8$ leaves $A_7 = SU(8)$ (dim~63)---the largest surviving sub-algebra among the three D4 legs.
    \item \textbf{Smallest Kac label:} $m_2 = 3$ is the minimum among D4-leg Kac labels, giving the most constrained $U(1)$ charge spectrum ($q_Y \in \{0, \pm 1, \pm 2, \pm 3\}$).
\end{enumerate}
The $U(1)_Y$ charge spectrum of the 120 positive $E_8$ roots decomposes as $28 + 56 + 28 + 8$ for charges $q_Y = 0, 1, 2, 3$ respectively.

\textbf{Classification:} The full chain 600-cell $\to$ 24-cell $\to$ $D_4$ $\to$ $E_8$ Dynkin $\to$ $SU(3) \times SU(2) \times U(1)$ is \textbf{DERIVED} from geometric and algebraic structure, with no free parameters. The identification of the 96-vertex complement as the matter sector is a PATTERN (the correct counting but no spinor representation derivation).

\subsection{Interaction Selection Rules and Propagator (EXP-136, 137)}
\label{sec:interactions}

The edge structure of the 600-cell imposes \textbf{automatic selection rules} on particle interactions. Since no edges connect gauge vertices to each other:
\begin{equation}
\text{A-A} = \text{A-B} = \text{B-B} = 0 \quad \text{(edges)}
\end{equation}
the only non-vanishing edges are A-C (96), B-C (192), and C-C (432), where C denotes the 96 fermion vertices. This propagates to interaction vertices: of the 10 possible triangle types and 15 possible tetrahedron types, only \textbf{three} survive in each case:

\begin{table}[H]
\centering
\begin{tabular}{lcccl}
\toprule
\textbf{Triangles} & \textbf{Count} & \textbf{Per fermion} & \textbf{Ratio} & \textbf{Physical analogue} \\
\midrule
ACC (gauge--fermion--fermion) & 240 & 5 & 1 & QED/QCD vertex \\
BCC (Higgs--fermion--fermion) & 480 & 10 & 2 & Yukawa coupling \\
CCC (fermion--fermion--fermion) & 480 & 15 & 3 & Fermion self-coupling \\
\midrule
\textbf{Total} & \textbf{1200} & \textbf{30} & & \\
\bottomrule
\end{tabular}
\caption{Triangle (3-point vertex) classification. Per-fermion counts are uniform across all 96 Type-C vertices. The 1:2:3 ratio is exact.}
\end{table}

Similarly, tetrahedra (4-point vertices) decompose as ACCC\,(160) : BCCC\,(320) : CCCC\,(120), with ratio 2:4:1.5 reflecting the same hierarchy.

\textbf{Key consequence:} Gauge bosons interact \emph{only} through fermion mediators. The ``effective gauge adjacency'' (2-step paths through fermions) connects every gauge pair by exactly 3 paths, with eigenvalues $\{24, 12, 0, -6, -12\}$---all integers despite the irrational ($\phi$-dependent) spectrum of the full graph.

\textbf{Propagator structure.} The massless Green's function $G = (L + 0)^{-1}_{\perp}$ (pseudoinverse of the Laplacian, projecting out the constant mode) shows \textbf{oscillatory behavior}:

\begin{center}
\begin{tabular}{crl}
$d$ & $\langle G_{\text{CC}}(d)\rangle$ & Sign \\
\hline
1 & $+0.01527$ & $+$ \\
2 & $+0.00236$ & $+$ \\
3 & $-0.00335$ & $-$ \\
4 & $-0.00634$ & $-$ \\
5 & $-0.00791$ & $-$ \\
\end{tabular}
\end{center}

The sign change at $d = 3$ is characteristic of \textbf{confinement} on finite graphs: beyond the interaction range of common gauge neighbors (which drops to zero at $d \geq 3$), the propagator becomes repulsive. At $d = 1$, the first eigenspace (Laplacian eigenvalue $\lambda_1 = 12 - 6\phi \approx 2.29$) contributes 77\% of the total amplitude, confirming that the lowest spectral mode dominates short-range interactions. The propagator values live in $\mathbb{Q}(\sqrt{5})$; for example, $G(1) = (-4 + 23\phi)/2175$.

\textbf{Classification:} Selection rules and triangle/tetrahedron counts are \textbf{DERIVED}. The mapping to SM interaction types (ACC $\to$ gauge vertex, BCC $\to$ Yukawa, CCC $\to$ self-coupling) is \textbf{PATTERN}. The confinement interpretation is \textbf{SPECULATIV}.

\subsection{The $1 + 3 + 8$ Edge Split and Effective 24-Cell (EXP-143, 144)}
\label{sec:edge_split}

Each fermion vertex has 12 neighbors: 1 of type A, 2 of type B, and 9 of type C. The Standard Model gauge group has dimension $1 + 3 + 8 = 12$, matching the total but not the internal structure ($1 + 2 + 9$ vs.\ $1 + 3 + 8$). A detailed analysis of the local neighborhood reveals a \textbf{natural geometric splitting} that resolves this discrepancy.

\textbf{Unique special CC edge.} Among the 9 C-neighbors of each fermion, exactly \textbf{one} has a distinct gauge-overlap profile:
\begin{itemize}
    \item Degree 4 in the local CC subgraph (vs.\ degree 3 for 6 others)
    \item Shares the same A-vertex (cross-polytope) gauge neighbor as the central fermion
    \item Shares \textbf{zero} B-vertex (tesseract) gauge neighbors
    \item Has minimum total gauge overlap (1 vs.\ 2 for the others)
\end{itemize}
This classification is \textbf{uniform} across all 96 fermion vertices, with exactly 48 such ``special'' CC edges in total. The per-fermion edge profile is:

\begin{center}
\begin{tabular}{ccccc}
\toprule
\textbf{Edge type} & \textbf{Count} & \textbf{shared A} & \textbf{shared B} & \textbf{SM sector} \\
\midrule
AC & 1 & -- & -- & $U(1)_Y$ \\
BC & 2 & -- & -- & $SU(2)_L$ (charged) \\
CC$_{\text{special}}$ & 1 & 1 & 0 & $SU(2)_L$ (neutral) \\
CC$_{\text{regular}}$ & 8 & mixed & mixed & $SU(3)_C$ \\
\midrule
\textbf{Total} & \textbf{12} & & & $\mathbf{1 + 3 + 8}$ \\
\bottomrule
\end{tabular}
\end{center}

\textbf{Effective gauge adjacency = 24-cell.} The Gram matrix $G = M_{\text{GF}} M_{\text{GF}}^T$, where $M_{\text{GF}}$ is the $24 \times 96$ gauge--fermion incidence matrix, encodes the effective gauge--gauge coupling mediated by fermions. Its off-diagonal entries take only values 0 or 3 (shared fermion neighbors), so $G = 12\mathbf{I} + 3\,A_{\text{eff}}$ where $A_{\text{eff}}$ has entries in $\{0, 1\}$. Remarkably:
\begin{equation}
A_{\text{eff}} = A_{\text{24-cell}} \quad \text{(verified element-by-element)}
\end{equation}
The effective adjacency through shared fermions is \textbf{identical} to the nearest-neighbor adjacency of the standard 24-cell (D$_4$ root system). Furthermore, this arises from the full 600-cell adjacency $A$ via:
\begin{equation}
A_{\text{eff}} = \frac{1}{3}\, A^2\big|_{\text{gauge}} \,,
\end{equation}
where $A^2|_{\text{gauge}}$ restricts the square of the adjacency matrix to the 24 gauge vertices. The factor $1/3$ is exact and universal.

\textbf{Hop-invariant eigenspaces.} The Gram matrix spectrum is $\{36, 24, 12, 6, 0\}$ with multiplicities $\{1, 4, 9, 8, 2\}$---all eigenvalues are multiples of 6, and the first three multiplicities are perfect squares. Higher-hop matrices $G_n = M_{\text{GF}} A_{\text{CC}}^n M_{\text{GF}}^T$ share the \textbf{same eigenspaces} for all $n$, with the eigenvalue $6$ (multiplicity 8) remaining \textbf{hop-invariant}. The dominant eigenvalue follows $36 \cdot 9^n$, where 9 is the CC-subgraph degree.

\textbf{Plaquette structure (EXP-145).} The $1 + 3 + 8$ edge classification induces a decomposition of all 1200 triangular plaquettes into five types:
\begin{center}
\begin{tabular}{lcc}
\toprule
\textbf{Type} & \textbf{Count} & \textbf{Pure/Mixed} \\
\midrule
$SU(3)$-$SU(3)$-$SU(3)$ & 288 & Pure \\
$SU(2)$-$SU(2)$-$SU(3)$ & 480 & Mixed \\
$SU(2)$-$SU(3)$-$SU(3)$ & 192 & Mixed \\
$SU(3)$-$U(1)$-$U(1)$ & 192 & Mixed \\
$SU(2)$-$U(1)$-$U(1)$ & 48 & Mixed \\
\bottomrule
\end{tabular}
\end{center}
Only $SU(3)$ has pure plaquettes; $U(1)$ and $SU(2)$ have \textbf{zero}. This is the lattice analogue of the fact that only gluons carry color charge and exhibit cubic self-interaction---the photon and $W/Z$ bosons do not self-interact in the pure gauge sector. The 288 pure $SU(3)$ plaquettes satisfy $288 = 24 \times 12 = \operatorname{Tr}(G)$.

\textbf{Vertex-transitivity and coupling independence (EXP-146).} The $H_4$ symmetry group acts transitively on the 120 vertices, with a fundamental consequence for the spectral action: every eigenspace of the adjacency matrix projects \textbf{identically} onto all three gauge sectors:
\[
\chi_k(\text{U1}) = \chi_k(\text{SU2}) = \chi_k(\text{SU3}) \quad \forall\, k = 0, \ldots, 8,
\]
where $\chi_k(S) = |E_S|^{-1} \sum_{(i,j) \in E_S} P_k[i,j]$ is the average projection of eigenspace $k$ onto edges of sector $S$. Similarly, the vertex-type weights are proportional to $|A|:|B|:|C| = 8:16:96$ in every eigenspace without exception. Consequently, no Connes--Chamseddine spectral action $\operatorname{Tr}(f(D^2))$ can differentiate coupling constants across sectors---the edge-type structure is invisible to the spectrum.

This establishes a \textbf{separation theorem}: the 600-cell encodes two independent geometric properties that correspond to two independent physical structures:
\begin{enumerate}
\item \textbf{Lagrangian form} (from topology): the $1 + 3 + 8$ edge split, plaquette types, selection rules---determining \emph{what} interacts with \emph{what}.
\item \textbf{Coupling values} (from the spectrum): eigenvalue ratios $\theta_1/\theta_3 = 2\phi$ and the intersection number $a_1 = 5$---determining \emph{how strongly} they interact.
\end{enumerate}
This mirrors the Standard Model, where the gauge group structure and coupling constant values are also logically independent.

\textbf{Shared-neighbor asymmetry (EXP-147).} While the spectral projection is sector-blind, the \emph{topological} neighborhood structure distinguishes sectors. Each edge has exactly $a_1 = 5$ shared neighbors (distance-regular), but their type decomposition is \textbf{exact} (standard deviation zero within each edge type):
\begin{center}
\begin{tabular}{lcccc}
\toprule
\textbf{Edge type} & \textbf{Count} & \textbf{shared A} & \textbf{shared B} & \textbf{shared C} \\
\midrule
$U(1)$ (AC) & 96 & 0 & 0 & 5 \\
$SU(2)_{\text{ch}}$ (BC) & 192 & 0 & 0 & 5 \\
$SU(2)_{\text{n}}$ (CC$_*$) & 48 & 1 & 0 & 4 \\
$SU(3)$ sub-type 1 & 96 & 0 & 1 & 4 \\
$SU(3)$ sub-type 2 & 96 & 0 & 2 & 3 \\
$SU(3)$ sub-type 3 & 192 & 1 & 1 & 3 \\
\bottomrule
\end{tabular}
\end{center}
$U(1)$ and $SU(2)_{\text{ch}}$ edges share \textbf{only fermion} neighbors (zero gauge-boson access), while $SU(3)$ edges share 1.75 gauge-boson neighbors on average across three sub-types. This is the lattice realization of Abelian vs.\ non-Abelian character: the photon interacts only through fermion loops, while gluons have direct gauge-boson access.

\textbf{Per-fermion gauge budget.} Each fermion has a total gauge access of exactly 15 (uniform across all 96 fermions): 0 from $U(1)$ edges, 1 from $SU(2)$ edges (solely from the neutral CC$_*$), and 14 from $SU(3)$ edges.

\textbf{Topological decomposition of $10\phi$.} The total fermion-shared neighbors across all $U(1)$ edges is $96 \times 5 = 480$, while the total gauge-shared neighbors across all $SU(2)$ edges is $48 \times 1 = 48$ (only the neutral CC$_*$ edges contribute). Their ratio is:
\begin{equation}
\frac{\text{fermion-shared}_{U(1)}}{\text{gauge-shared}_{SU(2)}} = \frac{480}{48} = 10\,.
\end{equation}
Since $\alpha_s/\alpha = 10\phi$ (Section~\ref{sec:alpha_s}), the factor 10 now has a topological origin: it counts how many fermion-loop mediations in the $U(1)$ sector correspond to one gauge self-interaction in the $SU(2)$ sector. The factor $\phi$ remains spectral (from $\theta_1/\theta_3 = 2\phi$, where the factor 2 itself equals $n_{U(1)}/n_{\text{CC}_*} = 96/48$). This provides a dual interpretation of the strong-to-electromagnetic coupling ratio: spectral ($a_1 \cdot \theta_1/\theta_3$) and topological ($\text{fermion}_{U(1)}/\text{gauge}_{SU(2)} \cdot \phi$).

\textbf{$SU(3)$ sub-type decomposition (EXP-148).} The 384 $SU(3)$ edges decompose into three sub-types based on their shared-neighbor profile, each forming a distinct sub-graph on the 96 fermion vertices:
\begin{center}
\begin{tabular}{lccccc}
\toprule
\textbf{Sub-type} & \textbf{Edges} & \textbf{Degree} & \textbf{Components} & \textbf{Per fermion} & \textbf{CC$_*$ link} \\
\midrule
T1 (0A, 1B, 4C) & 96 & 2 & 24 & 2 & 100\% \\
T2 (0A, 2B, 3C) & 96 & 2 & 32 & 2 & 0\% \\
T3 (1A, 1B, 3C) & 192 & 4 & \textbf{8} & 4 & 50\% \\
\bottomrule
\end{tabular}
\end{center}
The sub-type T3 has exactly \textbf{8 connected components}, each containing 12 fermion vertices---matching $\dim(SU(3)) = 8$ and the vertex degree 12. The mechanism is \emph{derived}: each T3 edge shares exactly one Type-A neighbor (by its profile $(1A,1B,3C)$), and the 8 Type-A vertices partition the 192 T3 edges into 8 disjoint groups of 24, each forming a single connected component (EXP-149). Thus $\dim(SU(3)) = |A| = 8$ is a geometric identity, not a numerical coincidence. Similarly, T1 components are labeled by the 24 A--B pairs each edge accesses, and T2 components by the 32 B--B pairs. The three sub-type adjacency matrices do \textbf{not commute} ($[A_{T_i}, A_{T_j}] \neq 0$ for all pairs), encoding non-Abelian structure consistent with $SU(3)$.

The CC$_*$ connectivity reveals a hierarchy: T1 edges are \emph{all} connected to the $SU(2)$ neutral partner (linking the strong and weak sectors), T2 edges are \emph{completely disconnected} from it, and T3 edges are connected at 50\%. The 288 pure $SU(3)$ plaquettes decompose as $192\,(\text{T1-T2-T3}) + 64\,(\text{T3-T3-T3}) + 32\,(\text{T2-T2-T2})$; the dominant mixed type reflects inter-sub-type coupling.

\textbf{Classification:} The $1 + 3 + 8$ edge split, effective adjacency $A_{\text{eff}} = A_{\text{24-cell}}$, plaquette decomposition, separation theorem, shared-neighbor profiles, the topological decomposition of $10\phi$, and the $SU(3)$ sub-type structure including the A-vertex generation of 8 components are \textbf{DERIVED}. The SM sector assignments remain \textbf{PATTERN}.

%==============================================================
\section{CKM and PMNS Mixing Matrices}
%==============================================================

\subsection{The Phi-Tower Formula}

Mixing angles follow a geometric hierarchy based on powers of the golden ratio:

\begin{equation}
\boxed{\theta_{ij} = \arctan(\phi^{-n_{ij}})}
\end{equation}

\subsection{CKM Matrix (Quarks)}

\textbf{Integer exponents.} The leading-order CKM hierarchy follows the phi-tower formula with integer exponents motivated by algebraic dimensions:

\begin{table}[H]
\centering
\begin{tabular}{ccccc}
\toprule
\textbf{Angle} & \textbf{$n$} & \textbf{Calculated} & \textbf{Experimental} & \textbf{Error} \\
\midrule
$\theta_{12}$ (Cabibbo) & 3 & $13.28°$ & $13.04°$ & \textbf{1.8\%} \\
$\theta_{23}$ & 7 & $1.97°$ & $2.34°$ & 15.6\% \\
$\theta_{13}$ & 12 & $0.18°$ & $0.22°$ & 18.7\% \\
\bottomrule
\end{tabular}
\caption{CKM mixing angles from phi-tower with integer exponents.}
\end{table}

\textbf{Coxeter exponent origin.} The exponents $\{3, 7, 12\}$ are expressible through the $H_4$ Coxeter exponents $\{1, 11, 19, 29\}$:
\begin{equation}
3 = \frac{1+11}{4}, \qquad 7 = \frac{|1-29|}{4}, \qquad 12 = \frac{19+29}{4}
\end{equation}
These are also compatible with algebraic dimension assignments: $n = 3 = \dim(\text{Im}(\mathbb{H}))$, $n = 7 = \dim(\text{Im}(\mathbb{O}))$, $n = 12 = $ vertex degree.

\textbf{Half-integer refinement.} Allowing half-integer exponents, a scan yields $\sin\theta_{ij} = \phi^{-n_{ij}}$ with $n = (3,\, 13/2,\, 23/2)$, giving mean error $\textbf{6.0\%}$ across all 9 CKM elements. The Cabibbo angle $\theta_C = \arctan(\phi^{-3}) = 13.28°$ remains the single strongest geometric prediction (1.8\% error, 0 free parameters).

\textbf{$E_8$ cross-pairing (EXP-129).} The half-integer exponents also arise from \emph{cross-pairings} of $H_4$ and $H_4'$ Coxeter exponents:
\begin{equation}
n_{12} = \frac{1+11}{4} = 3, \qquad n_{23} = \frac{19+7}{4} = 6.50, \qquad n_{13} = \frac{29+17}{4} = 11.50
\end{equation}
where $\{1,11,19,29\}$ are $H_4$ exponents and $\{7,17\}$ are $H_4'$ exponents from the $E_8$ branching $E_8 \to H_4 \oplus H_4'$. The cross-pairing $n_{23} = 6.50$ matches the fitted value $6.65$ to $2.3\%$, and $n_{13} = 11.50$ matches $11.57$ to $0.6\%$. This provides a cleaner $E_8$-based structure for the half-integer exponents, though the specific pairing rules remain underived.

\textbf{Wolfenstein parametrization.} Setting $\lambda = \phi^{-3}$ in the standard Wolfenstein expansion and fitting a single parameter $A = V_{cb}/\lambda^2$ yields mean error \textbf{3.1\%} with 1 free parameter.

\textbf{Classification:} The Cabibbo angle is a robust PATTERN ($1.8\%$, 0 free parameters). The full CKM matrix is PATTERN+FIT ($3{-}6\%$ with 0--1 free parameters). A complete derivation from first principles is not yet achieved.

\subsection{PMNS Matrix (Neutrinos)}

With perturbations from the charged lepton sector ($\epsilon \approx m_\mu/m_\tau \approx 0.06$):

\begin{table}[H]
\centering
\begin{tabular}{cccccc}
\toprule
\textbf{Angle} & \textbf{$n$} & \textbf{GUT} & \textbf{+ Pert.} & \textbf{Exp.} & \textbf{Error} \\
\midrule
$\theta_{12}$ (solar) & 1 & $31.72°$ & $33.3°$ & $33.41°$ & \textbf{0.3\%} \\
$\theta_{13}$ (reactor) & 4 & $8.30°$ & $8.5°$ & $8.54°$ & \textbf{0.5\%} \\
$\theta_{23}$ (atmospheric) & 0 & $45.0°$ & $47.7°$ & $49.0°$ & 2.7\% \\
\bottomrule
\end{tabular}
\caption{PMNS mixing angles with perturbations.}
\end{table}

\textbf{Exponent origins:}
\begin{itemize}
    \item $n=0$: Identity ($\phi^0 = 1$, maximal mixing)
    \item $n=1$: Edge/radius ratio = $1/\phi$ (fundamental 600-cell ratio)
    \item $n=4$: $\dim(\mathbb{R}^4) = 4$ (space dimensionality)
\end{itemize}

\textbf{Key insight:} Quarks ``see'' algebraic structure (quaternions, octonions), while leptons ``see'' geometric structure (space, ratios).

%==============================================================
\section{Lepton Mass Hierarchy from Holonomy}
%==============================================================

\subsection{The Mass Ratio Pattern}

The lepton mass ratios follow a remarkable pattern based on powers of $\phi$:

\begin{table}[H]
\centering
\begin{tabular}{lcccc}
\toprule
\textbf{Ratio} & \textbf{$n$} & \textbf{$\phi^n$} & \textbf{Experimental} & \textbf{Error} \\
\midrule
$m_\mu/m_e$ & 11 & 199.0 & 206.77 & 3.8\% \\
$m_\tau/m_\mu$ & 6 & 17.94 & 16.82 & 6.7\% \\
$m_\tau/m_e$ & 17 & 3571.0 & 3477.1 & 2.7\% \\
\bottomrule
\end{tabular}
\caption{Lepton mass ratios as powers of $\phi$.}
\end{table}

\subsection{Derivation from 600-Cell Geometry}

The exponents 5, 6, 11, 17 are \textbf{derived} from the intrinsic geometry of the 600-cell:

\begin{equation}
\boxed{
\begin{aligned}
5 &= \text{diameter of the 600-cell graph} \\
6 &= \text{decagons per vertex} \\
11 &= 5 + 6 \quad \text{(holonomy additivity)} \\
17 &= 11 + 6 \quad \text{(second excitation level)}
\end{aligned}
}
\end{equation}

\subsection{The Hopf Fibration Connection}

The 600-cell, inscribed in $S^3$, inherits the \textbf{Hopf fibration} structure $S^3 \to S^2$:

\begin{itemize}
    \item 120 vertices $\to$ 30 fibers $\times$ 4 vertices each
    \item Edge angle $= \pi/5 = 36°$ \textbf{(exact)}
    \item Total phase on decagon $= 2\pi$
\end{itemize}

\textbf{The holonomy argument:} When a spinor field is parallel-transported along a path on $S^3$:
\begin{itemize}
    \item Traversing the \textbf{diameter} (5 steps) accumulates phase $\sim 5 \cdot \log\phi$
    \item Each \textbf{decagon} (6 per vertex) contributes phase $\sim 1 \cdot \log\phi$
    \item Total for muon: $(5 + 6) \cdot \log\phi = 11 \cdot \log\phi$
\end{itemize}

The mass formula becomes:
\begin{equation}
m = m_e \cdot \phi^n = m_e \cdot e^{n \log\phi}
\end{equation}

\subsection{The Binary Icosahedral Group}

The holonomy group of the 600-cell is related to the \textbf{binary icosahedral group} $2I$:
\begin{itemize}
    \item $|2I| = 120$ (same as number of vertices!)
    \item Irreducible representations involve powers of $\phi$
    \item Holonomy is ``quantized'' in units related to $\phi$
\end{itemize}

\textbf{Conclusion:} The exponents 5, 6, 11, 17 are not numerological coincidences---they are \textbf{topological invariants} of the 600-cell lattice structure.

\subsection{Electroweak Phase Corrections (EXP-097)}

The bare exponents $n=11,\,6,\,17$ yield errors of 3.8\%, 6.7\%, and 2.7\% respectively. We decompose the effective exponents into their geometric components:

\begin{equation}
n_\mu = d_{\text{eff}} + k_{\text{eff}}, \quad
n_{\tau/\mu} = k_{\text{eff}}, \quad
n_\tau = d_{\text{eff}} + 2\,k_{\text{eff}}
\end{equation}

where $d_{\text{eff}} = 5 + \delta_d$ (diameter contribution) and $k_{\text{eff}} = 6 + \delta_k$ (decagonal contribution). A systematic numerical search (EXP-097) identifies:

\begin{equation}
\boxed{
\delta_d = \sin^2\theta_W = \frac{6}{26}, \qquad
\delta_k = -\frac{1}{\phi^4}
}
\end{equation}

\textbf{Physical interpretation:}
\begin{itemize}
    \item The diameter path ``feels'' the electroweak mixing angle --- the same ratio $6/26$ already derived from vertex structure counts.
    \item The decagonal loop ``feels'' a suppression of $1/\phi^4$, reflecting the four-dimensional embedding of the geodesic.
\end{itemize}

The corrected exponents give:

\begin{table}[H]
\centering
\begin{tabular}{lcccc}
\toprule
\textbf{Ratio} & \textbf{$n_{\text{eff}}$} & \textbf{Calculated} & \textbf{Experimental} & \textbf{Error} \\
\midrule
$m_\mu/m_e$ & 11.085 & 205.95 & 206.77 & \textbf{0.28\%} \\
$m_\tau/m_e$ & 16.939 & 1772.4 & 1776.9 & \textbf{0.25\%} \\
$m_\tau/m_\mu$ & 5.854 & 16.73 & 16.82 & \textbf{0.53\%} \\
\bottomrule
\end{tabular}
\caption{Lepton mass ratios with electroweak phase corrections.}
\end{table}

\textbf{Status:} The corrections $\sin^2\theta_W$ and $1/\phi^4$ are quantities already present in the theory (Sections 3.3 and 2.1), not free parameters. This reduces the lepton mass errors from 3--7\% to below 0.5\%.

\subsection{Unified Fermion Mass Formula (EXP-098--100)}

The electroweak phase corrections generalize to all fermions through sector-dependent corrections. The unified mass formula reads:

\begin{equation}
\boxed{m_f = m_e \cdot \phi^{\,a \cdot d_{\text{eff}} \;+\; b \cdot k_{\text{eff}}}}
\end{equation}

where $(a,b) \in \mathbb{Z}^2$ are \textbf{topological quantum numbers} specific to each fermion, and:
\begin{equation}
d_{\text{eff}} = 5 + \delta_d, \qquad k_{\text{eff}} = 6 + \delta_k
\end{equation}

The corrections $\delta_d, \delta_k$ depend on the fermion sector and reflect the dominant interaction:

\begin{table}[H]
\centering
\begin{tabular}{lccc}
\toprule
\textbf{Sector} & $\delta_d$ & $\delta_k$ & \textbf{Interpretation} \\
\midrule
Leptons ($Q=-1$) & $+\sin^2\theta_W$ & $-1/\phi^4$ & Electroweak \\
Up-type ($Q=+\tfrac{2}{3}$) & $+2\alpha_s/\pi$ & $+\alpha_s$ & QCD (positive shift) \\
Down-type ($Q=-\tfrac{1}{3}$) & $-1/5$ & $-\alpha_s$ & QCD (negative shift) \\
\bottomrule
\end{tabular}
\caption{Sector-dependent corrections. All quantities derive from the theory (Sections~3.2, 3.3).}
\end{table}

\textbf{Key point:} No new parameters are introduced. The lepton corrections use $\sin^2\theta_W$ (Section~3.3) and $\phi^{-4}$ (geometric). The quark corrections use $\alpha_s = 1/(2\phi^3)$ (Section~3.2) and $1/5 = 1/\text{diameter}$ (Section~2.1). The down-type correction $\delta_k = -\alpha_s = -1/(2\phi^3)$ is the same relation already established in Eq.~(6).

The complete fermion mass spectrum:

\begin{table}[H]
\centering
\begin{tabular}{lcccccr}
\toprule
\textbf{Fermion} & $Q$ & $(a,b)$ & $n$ & $m_{\text{calc}}$ & $m_{\text{exp}}$ & \textbf{Error} \\
\midrule
Electron & $-1$ & $(0,0)$ & 0 & 0.511 MeV & 0.511 MeV & 0.00\% \\
Muon & $-1$ & $(1,1)$ & 11 & 106.0 MeV & 105.7 MeV & \textbf{0.28\%} \\
Tau & $-1$ & $(1,2)$ & 17 & 1772 MeV & 1777 MeV & \textbf{0.25\%} \\
\midrule
Up & $+\tfrac{2}{3}$ & $(3,-2)$ & 3 & 2.15 MeV & 2.16 MeV & \textbf{0.30\%} \\
Charm & $+\tfrac{2}{3}$ & $(2,1)$ & 16 & 1283 MeV & 1270 MeV & \textbf{1.03\%} \\
Top & $+\tfrac{2}{3}$ & $(4,1)$ & 26 & 169.6 GeV & 172.8 GeV & 1.82\% \\
\midrule
Down & $-\tfrac{1}{3}$ & $(1,0)$ & 5 & 5.15 MeV & 4.67 MeV & $10.2\%^\dagger\!^\ddagger$ \\
Strange & $-\tfrac{1}{3}$ & $(1,1)$ & 11 & 87.3 MeV & 93.4 MeV & 6.6\% \\
Bottom & $-\tfrac{1}{3}$ & $(-1,4)$ & 19 & 4191 MeV & 4180 MeV & \textbf{0.26\%} \\
\bottomrule
\end{tabular}
\caption{Complete fermion mass predictions. $n = 5a + 6b$ is the bare exponent. $^\dagger$Down quark experimental mass: $4.67^{+0.48}_{-0.17}$ MeV (PDG~\cite{pdg2024}). $^\ddagger$The prediction 5.15\,MeV coincides with the 1$\sigma$ upper bound ($4.67 + 0.48 = 5.15$\,MeV), placing it within $1\sigma$ of the central value (EXP-128).}
\end{table}

Seven of nine fermions are predicted within 2\% of experiment. The two outliers---Down ($10\%$) and Strange ($7\%$)---are the lightest quarks, where non-perturbative QCD effects and experimental mass uncertainties are largest. A comprehensive investigation (EXP-128) shows that the Down quark prediction (5.15\,MeV) coincides exactly with the PDG 1$\sigma$ upper bound, placing it within $1\sigma$ of the experimental value. The tension between Down and Strange is structural: both lie on the $a = 1$ line in the $(a,b)$ lattice and share the same $\delta_d$ correction, so improving one necessarily worsens the other.

\textbf{Structural patterns:}
\begin{enumerate}
\item \textbf{Strange--Muon duality:} Both share $(a,b) = (1,1)$ and $n=11$, differing only in their sector corrections. At the QCD scale $\mu \approx 1.17$ GeV, their running masses coincide.
\item \textbf{Topological complexity:} The average $|a|+|b|$ increases with generation: Gen~1 (2.0) $<$ Gen~2 (2.3) $<$ Gen~3 (4.3), suggesting heavier fermions require more complex paths on the 600-cell graph.
\item \textbf{Self-consistency:} All correction terms ($\sin^2\theta_W$, $\alpha_s$, $1/\phi^4$, $1/5$) were independently derived elsewhere in the theory.
\end{enumerate}

\subsection{Derivation of $(a,b)$ from Complexity Minimization}
\label{sec:ab_derivation}

The topological quantum numbers $(a,b)$ are \emph{not} free parameters. They are uniquely determined by the 600-cell geometry plus a single principle.

\textbf{Theorem.} For each fermion with mass level $n$, the quantum numbers $(a,b)$ are the \emph{unique} solution minimizing topological complexity:
\begin{equation}
\boxed{(a,b) = \arg\min_{a_1 a + b_1 b = n} \bigl(|a| + |b|\bigr)}
\end{equation}
where $a_1 = 5$ and $b_1 = 6$ are the intersection numbers of the 600-cell (Section~\ref{sec:spectral}).

\textbf{Proof sketch.} The general integer solution of $5a + 6b = n$ is $a = a_0 + 6t$, $b = b_0 - 5t$ for $t \in \mathbb{Z}$. The complexity $f(t) = |a_0 + 6t| + |b_0 - 5t|$ is a convex piecewise-linear function with kinks at $t = -a_0/6$ and $t = b_0/5$. For $n \neq 0$ these kinks are at distinct positions, so $f(t)$ has a unique integer minimizer. $\square$

This has been verified computationally for all 9 fermions:

\begin{table}[H]
\centering
\begin{tabular}{lcccccc}
\toprule
\textbf{Fermion} & $n$ & \textbf{Predicted} $(a,b)$ & $|a|+|b|$ & \textbf{Unique?} & \textbf{Gap} \\
\midrule
Electron & 0 & $(0,0)$ & 0 & Yes & 11 \\
Down & 5 & $(1,0)$ & 1 & Yes & 9 \\
Muon/Strange & 11 & $(1,1)$ & 2 & Yes & 9 \\
Charm & 16 & $(2,1)$ & 3 & Yes & 7 \\
Tau & 17 & $(1,2)$ & 3 & Yes & 7 \\
Up & 3 & $(3,-2)$ & 5 & Yes & 1 \\
Bottom & 19 & $(-1,4)$ & 5 & Yes & 1 \\
Top & 26 & $(4,1)$ & 5 & Yes & 3 \\
\bottomrule
\end{tabular}
\caption{Complexity minimization predicts all 9 fermion $(a,b)$ uniquely. The ``Gap'' column shows $|a|+|b|$ of the next-best decomposition minus the minimum.}
\label{tab:complexity}
\end{table}

The result is robust: $L^1$, $L^2$, and $L^\infty$ norms all select the same $(a,b)$ for all 9 fermions. This reduces the free parameter count from 18 (nine fermions $\times$ two quantum numbers) to 9 (the mass levels $n$ alone).

\textbf{Status:} The quantum numbers $(a,b)$ are DERIVED from complexity minimization with intersection numbers $a_1 = 5$, $b_1 = 6$. The sector-dependent corrections are DERIVED from quantities already in the theory. The remaining input is the set of 9 mass levels $n$, which come from the DSI mechanism (Section~\ref{sec:dsi}) and experimental masses. The Down quark ($10\%$) remains an open problem.

%==============================================================
\section{Discrete Scale Invariance: The Mass Mechanism}
\label{sec:dsi}
%==============================================================

The unified mass formula of Section~7.6 raises a natural question: \emph{why} do fermion masses scale as powers of $\phi$? In this section we show that the answer is \textbf{Discrete Scale Invariance} (DSI) \cite{sornette1998}---the mass spectrum arises from a potential with periodicity in logarithmic energy space, and the 600-cell lattice provides both the scaling factor and the selection rules.

\subsection{The DSI Potential}

A scalar field on the 600-cell lattice experiences a self-similar potential with minima at discrete mass scales. The natural form is:

\begin{equation}
\boxed{V(\psi) = A \sin^2\!\left(\frac{\pi \ln(|\psi|/m_e)}{\ln\phi}\right)}
\label{eq:dsi_potential}
\end{equation}

where $A$ sets the barrier height. This potential has degenerate minima at $|\psi| = m_e \cdot \phi^n$ for all integers $n$, defining a \textbf{DSI ladder} of allowed mass levels.

\textbf{Why $\phi$?} The scaling factor is not a free parameter---it is fixed by the 600-cell geometry. The edge angle of the 600-cell inscribed in $S^3$ is $\pi/5$ (exact), and $2\cos(\pi/5) = \phi$. Under parallel transport, each edge contributes a holonomy phase $\ln\phi$ to a spinor field, so a path of $n$ edges yields a mass ratio $e^{n\ln\phi} = \phi^n$.

\subsection{Holonomy as Discrete Scale Invariance}

The equivalence between holonomy on $S^3$ and DSI is exact:

\begin{equation}
\text{Phase per edge} = \ln\phi \quad \Longrightarrow \quad \text{Mass ratio after } n \text{ edges} = \phi^n
\end{equation}

The Hopf fibration $S^3 \to S^2$ decomposes paths on the 600-cell into two fundamental classes:
\begin{itemize}
    \item \textbf{Diameter traversals:} 5 edges across the graph (graph diameter = 5)
    \item \textbf{Decagonal loops:} 6 edges around a decagonal geodesic (6 decagons per vertex)
\end{itemize}

A general path wrapping $a$ times across the diameter and $b$ times around decagonal loops traverses $n = 5a + 6b$ edges, giving:
\begin{equation}
m = m_e \cdot \phi^{5a + 6b}
\end{equation}

The coefficients 5 and 6 have a graph-theoretic origin: they are the \textbf{intersection numbers} $a_1 = 5$ and $b_1 = 6$ of the 600-cell association scheme (Section~\ref{sec:spectral}). Thus $n = a_1 a + b_1 b$, and the mass formula is fully determined by topological invariants.

With sector-dependent corrections (Section~7.5--7.6), this becomes the unified mass formula $m_f = m_e \cdot \phi^{a \cdot d_{\text{eff}} + b \cdot k_{\text{eff}}}$.

\textbf{Physical interpretation:} Fermions are \textbf{topological solitons} sitting at minima of the DSI potential. Each soliton is characterized by its winding numbers $(a,b)$ on the coarse 600-cell graph. The mass is not fitted---it is determined by which potential minimum the soliton occupies. The specific $(a,b)$ for each fermion follows from complexity minimization (Section~\ref{sec:ab_derivation}).

\subsection{Level Selection and Empty Levels}
\label{sec:level_selection}

The DSI ladder has minima at $m_e \cdot \phi^n$ for \emph{all} integers $n = 0, 1, 2, \ldots$. Since $\gcd(5,6)=1$, every non-negative integer $n$ can be written as $5a + 6b$ for some $(a,b) \in \mathbb{Z}^2$. Of the 27 levels from $n=0$ to $n=26$, only 8 are occupied by known fermions (with muon and strange sharing $n=11$).

Two algebraic constraints substantially narrow the allowed set:

\textbf{Condition 1: $\mathbb{Z}[\phi]$ norm constraint.} For each level $n$ with minimal-complexity representative $(a,b)$, compute the algebraic norm $N(a + b\phi) = a^2 + ab - b^2$. Occupied levels have $|N| \in \{0, 1, 5, 19\}$ only---excluding all levels with $|N| \in \{4, 9, 11, 16, 25, \ldots\}$. This reduces 27 candidates to 14.

\textbf{Condition 2: Three-line rule in the $(a,b)$ lattice.} The remaining selection is geometric: an allowed level is occupied only if its $(a,b)$ representative lies on
\begin{equation}
\{\text{origin}\} \;\cup\; \{a = 1\} \;\cup\; \{b = 1\} \;\cup\; \{3a + 2b = 5\}
\end{equation}

The three lines have a remarkable structure:
\begin{itemize}
    \item Line $a = 1$: $d(1,0)$, $\mu/s(1,1)$, $\tau(1,2)$ --- leptons + down quark
    \item Line $b = 1$: $\mu/s(1,1)$, $c(2,1)$, $t(4,1)$ --- heavy quarks
    \item Line $3a+2b = 5$: $u(3,-2)$, $\mu/s(1,1)$, $b(-1,4)$ --- cross-generation
\end{itemize}

All three lines intersect at $(1,1)$ = muon/strange, the triple intersection point. This explains the $\mu$--$s$ degeneracy ($n=11$ shared) as a topological consequence: $(1,1)$ is the unique junction of all fermion lines.

Together, the two conditions correctly classify \textbf{28 of 31 levels} ($n = 0$ to $30$). Three false positives remain: $n = 1\,(-1,1)$, $n = 6\,(0,1)$, $n = 23\,(1,3)$---all on lines $a=1$ or $b=1$ but unoccupied.

\textbf{Per-line refinement (EXP-123).} The three false positives are eliminated by per-line constraints, one of which is derivable:
\begin{itemize}
    \item \textbf{Line $b = 1$:} Require the \emph{signed} norm $N(a + b\phi) = a^2 + ab - b^2 > 0$. On this line, $N = a^2 + a - 1$, which is negative for $a \leq 0$ and positive for $a \geq 1$. This excludes both false positives $n = 1\,(a=-1)$ and $n = 6\,(a=0)$. The occupied values $a = \{1, 2, 4\} = \{2^0, 2^1, 2^2\}$ give norms $|N| = \{1, 5, 19\}$---exactly the three non-zero allowed magnitudes. \textbf{Status: DERIVED} (norm positivity in $\mathbb{Z}[\phi]$).
    \item \textbf{Line $a = 1$:} Require $b \leq 2$, excluding $n = 23\,(b=3)$. This is equivalent to limiting the number of generations to three. \textbf{Status: PATTERN} (physically motivated but not geometrically derived).
    \item \textbf{Line $3a + 2b = 5$:} No additional constraint needed; all three solutions are occupied ($3/3$). \textbf{Status: DERIVED} (self-consistent).
\end{itemize}

With these refinements, the selection rule achieves \textbf{31/31} accuracy. The three-generation limit ($b \leq 2$) remains the sole underived component.

\textbf{Classification:} The $\mathbb{Z}[\phi]$ norm constraint is DERIVED. The three-line rule is a PATTERN. The $b = 1$ line refinement via norm positivity is DERIVED. The $a = 1$ line refinement ($b \leq 2 =$ three generations) is PATTERN---deriving the number of fermion generations from geometry remains one of the deepest open problems in theoretical physics (EXP-127).

\subsection{Fractal Self-Similarity}

The symmetry group of the 600-cell satisfies a suggestive identity:
\begin{equation}
|H_4| = 14{,}400 = 120^2
\end{equation}

If interpreted as a \textbf{fractal hierarchy}---each of the 120 vertices of a coarse 600-cell resolving into a finer 600-cell of 120 vertices---then:
\begin{itemize}
    \item \textbf{Level 0} (coarse graph): determines the coupling constants ($\alpha$, $\alpha_s$, $\sin^2\theta_W$) from spectral properties
    \item \textbf{Level 1} (fine structure): determines fermion masses from soliton paths $(a,b)$ on the coarse graph
    \item Sector corrections ($\delta_d$, $\delta_k$) encode the coupling between levels
\end{itemize}

The number of DSI levels from the Planck scale to the electron mass is $\ln(m_P/m_e)/\ln\phi \approx 107$, and the spectral equation (Eq.~3) provides the self-consistency condition linking the lattice scaling factor $\phi$ to the coupling constant $\alpha$.

\textbf{Classification:} The DSI potential (Eq.~\ref{eq:dsi_potential}) and the scaling factor $\phi$ are DERIVED from 600-cell geometry. The mass formula $n = a_1 a + b_1 b$ with $a_1 = 5$, $b_1 = 6$ is DERIVED from intersection numbers. The quantum numbers $(a,b)$ are DERIVED from complexity minimization (Section~\ref{sec:ab_derivation}). The fractal interpretation ($|H_4| = 120^2$) is an OBSERVATION whose physical significance is SPECULATIVE. The occupied level set $\{0, 3, 5, 11, 16, 17, 19, 26\}$ remains PATTERN.

%==============================================================
\section{E8 Extension and Gauge Unification}
%==============================================================

\subsection{The 600-Cell to E8 Connection via the Icosian Lattice}
\label{sec:e8}

The 240 roots of $E_8$ are constructed explicitly from the 600-cell using the \textbf{icosian lattice} of Conway and Sloane \cite{conway1998icosian}.

\textbf{Step 1: Decomposition.} Each 600-cell vertex $v \in \mathbb{R}^4$ is decomposed as $v_i = a_i + b_i\phi$ with $a_i, b_i \in \frac{1}{2}\mathbb{Z}$, yielding an 8-tuple $(a_0, b_0, a_1, b_1, a_2, b_2, a_3, b_3) \in \mathbb{R}^8$.

\textbf{Step 2: Icosian norm.} The icosian inner product uses the Gram matrix (per coordinate pair):
\begin{equation}
G = \begin{pmatrix} 1 & 1 \\ 1 & 2 \end{pmatrix}, \qquad \det G = 1 \text{ (unimodular)}
\end{equation}
The icosian norm is $N_{\text{icos}}(v) = \sum_{i=0}^{3}(a_i^2 + 2a_i b_i + 2b_i^2)$. All 120 unit 600-cell vertices have $N_{\text{icos}} = 1$.

\textbf{Step 3: Cholesky embedding.} The Cholesky factorization $G = LL^T$ with $L = \bigl(\begin{smallmatrix}1&0\\1&1\end{smallmatrix}\bigr)$ gives the map $(a_i, b_i) \mapsto (a_i + b_i,\, b_i)$. Under this map, icosian norm equals Euclidean norm$^2$. The 120 vertices embed as a set $S \subset \mathbb{R}^8$ with $\|s\|^2 = 1$.

\textbf{Step 4: Second 600-cell.} The Galois conjugate $\phi' = (1-\sqrt{5})/2$ defines a second copy: $T = \phi' \cdot S$. Since $\phi\phi' = -1$, the conjugation acts on decomposition coefficients as $(a,b) \mapsto (a-b,\,-a)$. The icosian norm of $\phi'$ is $N_{\text{icos}}(\phi') = 1$, so $T$ also consists of unit icosians.

\textbf{Step 5: E8 root system.}
\begin{equation}
\boxed{S \cup T = 240 \text{ vectors}, \quad S \cap T = \emptyset}
\end{equation}

After rescaling by $\sqrt{2}$ to match the $E_8$ convention $\|r\|^2 = 2$, the inner product distribution is:
\begin{equation}
\langle r_i, r_j \rangle \in \{-2, -1, 0, +1\}, \quad \text{with multiplicities } 120,\; 6720,\; 15120,\; 6720
\end{equation}
which is the \textbf{exact} $E_8$ root system inner product structure. A Procrustes rotation matches all 240 vectors to the standard $E_8$ basis (240/240 exact).

\textbf{Classification:} DERIVED. The $E_8$ root system is an exact consequence of the 600-cell geometry via the icosian lattice construction. No free parameters.

\textbf{Physical interpretation:}
\begin{itemize}
    \item Low energies ($\sim M_Z$): spacetime ``sees'' only $S$ (the 600-cell)
    \item High energies ($\sim M_{\text{GUT}}$): spacetime ``sees'' $S \cup T$ (full $E_8$)
    \item The Galois conjugation $\phi \to \phi'$ acts as a ``mirror'' symmetry connecting the two copies
\end{itemize}

\subsection{The 120-Cell Dual: Spectral Structure and Fibonacci Discriminants}
\label{sec:120cell}

The 120-cell $\{5,3,3\}$ is the dual of the 600-cell $\{3,3,5\}$: its 600 vertices correspond to the 600 tetrahedral cells, with two vertices adjacent when their tetrahedra share a face. A universal \textbf{factor of 3} governs the duality:

\begin{table}[H]
\centering
\begin{tabular}{lccc}
\toprule
\textbf{Property} & \textbf{600-cell} & \textbf{120-cell} & \textbf{Ratio} \\
\midrule
Vertex degree & 12 & 4 & 3 \\
Graph diameter & 5 & 15 & 3 \\
Distinct eigenvalues & 9 & 27 & 3 \\
\bottomrule
\end{tabular}
\caption{The factor-3 duality between the 600-cell and 120-cell.}
\end{table}

The spectral structure of the 120-cell reveals a fundamental distinction. While the 600-cell spectrum is \textbf{spectrally pure}---all eigenvalues lie in $\mathbb{Q}(\sqrt{5})$ and all multiplicities are perfect squares---the 120-cell spectrum involves \emph{five} algebraic number fields:

\begin{table}[H]
\centering
\begin{tabular}{lccl}
\toprule
\textbf{Field} & \textbf{Eigenvalues} & \textbf{Total dim.} & \textbf{Discriminant} \\
\midrule
$\mathbb{Q}$ & 5 & 75 & -- \\
$\mathbb{Q}(\sqrt{5})$ & 10 & 182 & $5 = F_5$ \\
$\mathbb{Q}(\sqrt{2})$ & 2 & 96 & $2 = F_3$ \\
$\mathbb{Q}(\sqrt{13})$ & 2 & 32 & $13 = F_7$ \\
$\mathbb{Q}(\sqrt{21})$ & 2 & 32 & $21 = F_8$ \\
Cubic & 6 & 183 & -- \\
\midrule
Total & 27 & 600 & \\
\bottomrule
\end{tabular}
\caption{Algebraic field decomposition of the 120-cell spectrum. All quadratic discriminants are Fibonacci numbers $F_n$.}
\label{tab:120cell_spectrum}
\end{table}

The \textbf{Fibonacci discriminant pattern} is striking: every quadratic extension appearing in the 120-cell spectrum involves $\sqrt{F_n}$ for Fibonacci numbers $F_3 = 2$, $F_5 = 5$, $F_7 = 13$, $F_8 = 21$. The 600-cell, being spectrally pure in $\mathbb{Q}(\sqrt{5})$, uses only $\sqrt{F_5}$.

Two further structural features emerge:
\begin{enumerate}
\item \textbf{First-six matching:} The six largest eigenvalues of the 120-cell have multiplicities $1^2, 2^2, 3^2, 4^2, 5^2, 6^2$---identical to the 600-cell pattern. The departure occurs only at the 7th eigenvalue.
\item \textbf{Twin cubics:} Two irreducible cubic polynomials $x^3 - x^2 - 7x + 4$ (mult.~$25 = 5^2$) and $x^3 - x^2 - 7x + 8$ (mult.~$36 = 6^2$) share the same $x^2$- and $x$-coefficients, differing only in their constant terms (4 and 8). Galois conjugate eigenvalues always carry the same multiplicity.
\end{enumerate}

The intersection numbers of the 120-cell are $a_1 = 0$, $b_1 = 3$ (cf.\ $a_1 = 5$, $b_1 = 6$ for the 600-cell), and it is \emph{not} fully distance-regular (failing at $d = 4$). This spectral impurity supports the interpretation that the 600-cell is the fundamental geometric object, with the 120-cell playing a derived role.

\subsection{Unification Results}

Using MSSM beta functions as an $E_8$ proxy:

\begin{table}[H]
\centering
\begin{tabular}{lcc}
\toprule
\textbf{Model} & \textbf{Unification Scale} & \textbf{Discrepancy} \\
\midrule
Standard Model & -- & 27.9\% \\
600-cell + $E_8$ & $\sim 3 \times 10^{17}$ GeV & \textbf{2.4\%} \\
\bottomrule
\end{tabular}
\caption{Gauge unification comparison.}
\end{table}

%==============================================================
\section{Lagrangian Validation}
%==============================================================

\subsection{The Spectral Action}

Following the Spectral Action Principle of Chamseddine and Connes \cite{chamseddine1997}, the total action is defined through the Dirac operator $\mathcal{D}$ on the 600-cell lattice:

\begin{equation}
S = \text{Tr}(f(\mathcal{D}/\Lambda))
\end{equation}

where $\Lambda \sim M_{\text{Planck}}$ is the cutoff scale and $f$ is a cutoff function.

\subsection{Standard Tests (EXP-086)}

\begin{table}[H]
\centering
\begin{tabular}{lcc}
\toprule
\textbf{Test} & \textbf{Result} & \textbf{Error} \\
\midrule
Euler-Lagrange equations & Consistent & -- \\
Lorentz force recovery & Verified & -- \\
Higgs mass & 125.36 GeV & 0.09\% \\
Higgs VEV & 246.2 GeV & 0.0\% \\
Linear dispersion ($n=1$) & \textbf{Excluded} & -- \\
Quadratic dispersion ($n=2$) & Consistent & -- \\
\bottomrule
\end{tabular}
\caption{Standard Lagrangian validation tests. Gravitational tests (light deflection, precession) are presented in Appendix~A.}
\end{table}

\subsection{Advanced Tests (EXP-087)}

\subsubsection{Trans-Planckian Unitarity}

Unitarity is ensured by:
\begin{itemize}
    \item Hermitian Hamiltonian (real Laplacian spectrum)
    \item Finite derivatives on discrete lattice (no Ostrogradsky ghosts)
    \item $H_4$ symmetry (cancellation of negative-norm states)
\end{itemize}

\subsubsection{Weyl Conformal Invariance}

On the 600-cell (approximated as $S^3$):
\begin{itemize}
    \item $\chi = 0 \Rightarrow a \cdot E_4 = 0$ (Euler term vanishes)
    \item Constant curvature $\Rightarrow C_{\mu\nu\rho\sigma} = 0$ (Weyl tensor vanishes)
    \item Conformal anomaly is \textbf{zero} at leading order
\end{itemize}

\subsubsection{Cosmological Constant}

The Heat Kernel expansion suggests a suppression mechanism:

\begin{equation}
\Lambda_{\text{eff}} \sim \Lambda_{\text{natural}} \times \left(\frac{l_P}{L_{\text{Hubble}}}\right)^2 \sim 10^{-122} \times \Lambda_{\text{natural}}
\end{equation}

This yields the correct order of magnitude for the observed cosmological constant.

%==============================================================
\section{Experimental Predictions}
%==============================================================

\subsection{Photon Dispersion from GRBs}

The discrete spacetime structure predicts energy-dependent photon velocities:

\begin{equation}
\Delta t = \xi \left(\frac{E}{E_{\text{Planck}}}\right)^n \frac{d}{c}
\end{equation}

\textbf{Linear model ($n=1$):} \textbf{EXCLUDED}
\begin{itemize}
    \item Prediction: $\Delta t \sim 8$ s for 1 TeV photon at 1 Gpc
    \item Limit from GRB 090510 \cite{fermi2009}: $< 1$ ms
\end{itemize}

\textbf{Quadratic model ($n=2$):} CONSISTENT
\begin{itemize}
    \item Prediction: $\Delta t \sim 10^{-12}$ ms for 1 TeV at 1 Gpc
    \item Testable with CTA (Cherenkov Telescope Array) at higher energies
\end{itemize}

\subsection{Other Predictions}

\begin{enumerate}
    \item Running of coupling constants should follow 600-cell constraints
    \item \textbf{Sterile neutrino sector:} The 4-to-3 generation mechanism predicts a fourth group of fermions that are disconnected from the temporal axis. These are candidates for:
    \begin{itemize}
        \item Right-handed (sterile) neutrinos
        \item keV-scale dark matter particles
        \item Leptogenesis mediators
    \end{itemize}
    \item \textbf{Quark mass hierarchy:} The up quark mass follows $m_u/m_e = \phi^3$ with 0.2\% error, where $n=3 = \dim(\text{Im}(\mathbb{H}))$. The strange quark shares the muon exponent $n=11$, and at the scale $\mu \approx 1.17$ GeV, $m_s(\mu) = m_\mu$ exactly (3.5\% via 1-loop QCD running).
    \item CP-violating phase $\delta$ should have geometric origin
\end{enumerate}

%==============================================================
\section{Summary of Results}
%==============================================================

\begin{table}[H]
\centering
\begin{tabular}{lccc}
\toprule
\textbf{Quantity} & \textbf{Formula} & \textbf{Error} & \textbf{Status} \\
\midrule
\multicolumn{4}{l}{\textit{Coupling constants:}} \\
$\alpha$ & Spectral equation & 0.0001\% & Derived \\
$\alpha_s$ & $1/(2\phi^3)$ & 0.03\% & Derived \\
$\alpha_s/\alpha$ & $10\phi$ & Exact & Derived \\
$\sin^2\theta_W$ & $6/26$ & 0.19\% & Derived \\
$m_H$ & $m_W(\phi - 8\alpha)$ & 0.09\% & Derived \\
$\alpha_G$ & $\alpha^{8\phi^2}$ & 0.5\% & Derived \\
$m_e/m_P$ & $\alpha^{4\phi^2}$ & 0.16\% & Derived \\
\midrule
\multicolumn{4}{l}{\textit{Lepton mass hierarchy (from 600-cell holonomy):}} \\
$m_\mu/m_e$ & $\phi^{11}$ ($n=5+6$) & 3.8\% & \textbf{Derived} \\
$m_\mu/m_e$ (corrected) & $\phi^{5+\sin^2\theta_W+6-\phi^{-4}}$ & \textbf{0.28\%} & \textbf{Derived} \\
$m_\tau/m_\mu$ & $\phi^{6}$ & 6.7\% & \textbf{Derived} \\
$m_\tau/m_e$ (corrected) & $\phi^{d_{\text{eff}}+2k_{\text{eff}}}$ & \textbf{0.25\%} & \textbf{Derived} \\
$m_\tau/m_e$ & $\phi^{17}$ ($n=11+6$) & 2.7\% & \textbf{Derived} \\
$m_u/m_e$ & $\phi^{3}$ & 0.2\% & \textbf{Derived} \\
\midrule
\multicolumn{4}{l}{\textit{Unified fermion mass formula:}} \\
Full formula & $m_f = m_e \cdot \phi^{a \cdot d_{\text{eff}} + b \cdot k_{\text{eff}}}$ & 2.3\% avg & Derived \\
$m_c$ (charm) & $(2,1)$, up-type corr. & 1.03\% & Derived \\
$m_b$ (bottom) & $(-1,4)$, down-type corr. & 0.26\% & Derived \\
$m_t$ (top) & $(4,1)$, up-type corr. & 1.82\% & Derived \\
\midrule
\multicolumn{4}{l}{\textit{DSI mechanism (Section~\ref{sec:dsi}):}} \\
Scaling factor & $\phi$ from edge angle $\pi/5$ & Exact & Derived \\
DSI potential & $V \propto \sin^2(\pi\ln|\psi|/\ln\phi)$ & -- & Derived \\
Level selection & $n = 5a + 6b$ & -- & Derived \\
$(a,b)$ assignment & $\arg\min|a|+|b|$ s.t. $5a+6b=n$ & 9/9 & Derived \\
\midrule
\multicolumn{4}{l}{\textit{Mixing angles:}} \\
$\theta_C$ (CKM) & $\arctan(\phi^{-3})$ & 1.8\% & Pattern \\
CKM (Coxeter) & $n = H_4\text{ exps}/4$ & 6.0\% & Pattern \\
CKM ($E_8$ cross) & $n = (H_4 + H_4')/4$ & 0.6--2.3\% & Pattern \\
$\theta_{12}$ (PMNS) & $\arctan(1/\phi) + \epsilon$ & 0.3\% & Pattern+Fit \\
$\theta_{13}$ (PMNS) & $\arctan(\phi^{-4}) + \epsilon$ & 0.5\% & Pattern+Fit \\
$\theta_{23}$ (PMNS) & $45^\circ + \epsilon \cdot 45^\circ$ & 2.7\% & Pattern+Fit \\
\midrule
\multicolumn{4}{l}{\textit{E8 and level selection (Sections~\ref{sec:e8}, \ref{sec:level_selection}):}} \\
$E_8$ roots & $S \cup T$ (icosian) & 240/240 & \textbf{Derived} \\
Gauge group & 24-cell $\to D_4 \xrightarrow{E_8}$ SM (Kac $\{3,4,5\}$) & Exact & \textbf{Derived} \\
Level selection & $|N| \in \{0,1,5,19\}$ + 3 lines + $N>0$ & 31/31$^*$ & Pattern+Derived \\
\midrule
\multicolumn{4}{l}{\textit{Interaction structure (Section~\ref{sec:interactions}):}} \\
Selection rules & A-A $=$ A-B $=$ B-B $= 0$ & Exact & \textbf{Derived} \\
Triangle types & ACC : BCC : CCC $= 1:2:2$ & Exact & \textbf{Derived} \\
Per-fermion ratio & $5:10:15 = 1:2:3$ & Exact & \textbf{Derived} \\
Propagator sign & Oscillatory at $d \geq 3$ & -- & Derived \\
Edge split & $1 + 3 + 8 = \dim(\text{SM})$ & Exact & \textbf{Derived} \\
Eff.\ gauge adj. & $A_{\text{eff}} = A_{\text{24-cell}} = A^2|_g/3$ & Exact & \textbf{Derived} \\
Gram spectrum & $\{36,24,12,6,0\} \times \{1,4,9,8,2\}$ & Exact & \textbf{Derived} \\
Pure plaquettes & Only $SU(3)$: 288 of 1200 & Exact & \textbf{Derived} \\
Separation thm. & Form (topology) $\perp$ values (spectrum) & -- & \textbf{Derived} \\
Shared-nbr profiles & Exact per edge type (std $= 0$) & Exact & \textbf{Derived} \\
$10\phi$ decomp. & $480/48 \cdot \phi = 10\phi$ & Exact & \textbf{Derived} \\
$SU(3)$ sub-types & T3: 8 comp.\ $= |A| = \dim(SU(3))$ & Exact & \textbf{Derived} \\
Non-commutativity & $[A_{T_i}, A_{T_j}] \neq 0$ & Exact & \textbf{Derived} \\
A-vertex mechanism & 8 T3 comp.\ $=$ 8 A-vertices & Exact & \textbf{Derived} \\
Mixed-scale ratio & $\alpha_s(M_Z)/\alpha(0) = 10\phi$ & 0.03\% & Clarifying \\
\bottomrule
\end{tabular}
\caption{Summary of particle physics results. $^*$31/31 with per-line refinement; the $b \leq 2$ constraint (three generations) remains a pattern. Gravitational sector results are presented separately in Appendix~A (Table~\ref{tab:stg_vs_gr}).}
\end{table}

%==============================================================
\section{Conclusions}
\label{sec:conclusions}
%==============================================================

The 600-cell framework provides a coherent geometric basis for understanding the origin of fundamental constants in the particle physics sector. The key achievements are:

\begin{enumerate}
    \item \textbf{Numerical precision:} The coupling constants $\alpha$ (0.0001\%), $\alpha_s$ (0.03\%), $\sin^2\theta_W$ (0.19\%), and the Higgs mass (0.09\%) are derived from spectral and geometric properties of the 600-cell, not fitted.

    \item \textbf{Lepton mass hierarchy:} Exponents $n=11=5+6$ and $n=17=11+6$ are \textbf{derived} from the graph diameter (5) and decagons per vertex (6)---topological invariants of the 600-cell. Electroweak phase corrections ($\delta_d = \sin^2\theta_W$, $\delta_k = -1/\phi^4$) reduce errors to \textbf{0.25--0.53\%}.

    \item \textbf{Unified mass formula:} $m_f = m_e \cdot \phi^{a \cdot d_{\text{eff}} + b \cdot k_{\text{eff}}}$ predicts seven of nine fermion masses within 2\%. The sector-dependent corrections use only quantities already derived within the theory.

    \item \textbf{DSI mechanism:} The mass formula is not an empirical fit but arises from a Discrete Scale Invariance potential (Eq.~\ref{eq:dsi_potential}) whose scaling factor $\phi$ is fixed by the 600-cell geometry. Fermions are topological solitons whose mass levels are selected by their winding numbers $(a,b)$ on the lattice. Crucially, $(a,b)$ are not free parameters: they are uniquely determined by complexity minimization on the graph (Section~\ref{sec:ab_derivation}).

    \item \textbf{Spectral structure:} The 600-cell adjacency spectrum has 9 distinct eigenvalues, all of the form $a + b\phi$, with multiplicities that are all perfect squares $n^2$ ($n = 1, \ldots, 6$). The intersection numbers $a_1 = 5$, $b_1 = 6$ provide a graph-theoretic derivation of $\alpha_s/\alpha = 10\phi$ and a second independent path to $\sin^2\theta_W = 6/26$.

    \item \textbf{Gauge group derivation (complete):} The 24 non-fermionic vertices form the $D_4$ root system. The $E_8$ embedding breaks $D_4$ triality: the three legs extend with tail lengths $\{0, 1, 3\}$ and distinct Kac labels $\{3, 4, 5\}$, uniquely assigning each leg to a SM factor (Section~\ref{sec:gauge}). The final step---$SU(2)_R \to U(1)_Y$---follows from the dead-end isolation of the $m=3$ leg: with no non-Abelian extension to protect it, $SU(2)_R$ naturally reduces to its Cartan subalgebra $U(1)_Y$. Three independent structural arguments (topological isolation, maximal commutant $A_7$ with dim~63, smallest Kac label) confirm the breaking. The full chain is: 600-cell $\to$ 24-cell $\to$ $D_4$ $\to$ $E_8$ Dynkin $\to$ SM.

    \item \textbf{Fermion mapping:} Three generations encoded in $96 = 16 \times 3 \times 2$ vertices, with a geometric 4-to-3 selection mechanism via the temporal axis.

    \item \textbf{Mixing angles:} CKM exponents $\{3, 7, 12\}$ connect to $H_4$ Coxeter exponents. Cross-pairings of $H_4$ and $H_4'$ exponents from the $E_8$ branching provide a cleaner structure for the half-integer refinement (Section~6.2). PMNS from phi-tower with geometric exponent origins.

    \item \textbf{120-cell dual:} The dual polytope's spectrum spans five algebraic fields, with all quadratic discriminants being Fibonacci numbers ($F_3, F_5, F_7, F_8$). This spectral impurity confirms the 600-cell as the fundamental object.

    \item \textbf{E8 from 600-cell:} The icosian lattice construction (Section~\ref{sec:e8}) produces the 240 $E_8$ roots exactly from two copies of the 600-cell ($S$ and $T = \phi' \cdot S$), with zero overlap and exact inner product matching. This is a derivation, not a coincidence.

    \item \textbf{Level selection progress:} Two algebraic conditions---$\mathbb{Z}[\phi]$ norm restriction and a three-line rule in the $(a,b)$ lattice---plus norm positivity on the $b=1$ line achieve \textbf{31/31} classification accuracy (Section~\ref{sec:level_selection}). The sole remaining pattern is $b \leq 2$ on the $a=1$ line, equivalent to the three-generation limit. The triple intersection point $(1,1)$ explains the muon--strange degeneracy.

    \item \textbf{Gauge unification:} $E_8$ extension reduces discrepancy from 27.9\% to 2.4\%.

    \item \textbf{Validated Lagrangian:} Passes all tests---unitarity, Weyl invariance, correct Higgs sector.

    \item \textbf{Interaction selection rules:} The edge structure automatically forbids A-A, A-B, and B-B connections, leaving only three types of interaction vertex (ACC, BCC, CCC) in the exact ratio $1:2:2$. Each fermion participates in $5 + 10 + 15 = 30$ triangles with ratio $1:2:3$. The massless propagator changes sign at graph distance $d = 3$, exhibiting confinement-like behavior on the finite lattice (Section~\ref{sec:interactions}).

    \item \textbf{Edge split $\mathbf{1+3+8}$ (new):} Among the 9 C-neighbors of each fermion, exactly one has a geometrically unique gauge-overlap profile (degree 4 in local CC graph, shared A but zero shared B). Reclassifying this single edge yields $1\,(\text{AC}) + 3\,(2\,\text{BC} + 1\,\text{CC}_*) + 8\,\text{CC} = 12 = \dim(\text{SM})$. The effective gauge adjacency through shared fermions reproduces the 24-cell \emph{exactly}: $A_{\text{eff}} = A^2|_{\text{gauge}}/3$ (Section~\ref{sec:edge_split}).

    \item \textbf{Plaquette purity and separation theorem (new):} Among 1200 triangular plaquettes, only $SU(3)$ has pure plaquettes (288)---$U(1)$ and $SU(2)$ have zero---reproducing the non-Abelian self-interaction structure of QCD. A vertex-transitivity theorem proves that spectral projections are sector-blind ($\chi_k$ identical across U(1), SU(2), SU(3) for all 9 eigenspaces), establishing that the Lagrangian \emph{form} (from topology) and coupling \emph{values} (from the spectrum) are two independent geometric properties of the 600-cell (Section~\ref{sec:edge_split}).

    \item \textbf{Shared-neighbor asymmetry and topological $10\phi$ (new):} Each edge type has an \emph{exact} shared-neighbor profile (standard deviation zero): $U(1)$ and $SU(2)_{\text{ch}}$ edges share only fermion (C-type) neighbors, while $SU(3)$ edges share 1.75 gauge-boson neighbors on average---the topological origin of Abelian vs.\ non-Abelian character. The ratio $\alpha_s/\alpha = 10\phi$ decomposes as $(\text{fermion-shared}_{U(1)} / \text{gauge-shared}_{SU(2)}) \cdot \phi = (480/48) \cdot \phi$, giving the factor 10 a topological interpretation independent of the spectral derivation (Section~\ref{sec:edge_split}).

    \item \textbf{$SU(3)$ sub-type decomposition and A-vertex mechanism (new):} The 384 $SU(3)$ edges split into three sub-types by shared-neighbor profile: $T_1(0A,1B,4C) = 96$ edges with 24 components, $T_2(0A,2B,3C) = 96$ edges with 32 components, and $T_3(1A,1B,3C) = 192$ edges with \textbf{8 connected components}. The mechanism is now \emph{derived}: each T3 edge shares exactly one Type-A vertex, and the 8 A-vertices partition T3 into 8 disjoint connected groups of 24 edges, giving $\dim(SU(3)) = |A| = 8$ as a geometric identity (EXP-149). The sub-type adjacency matrices do not commute ($[A_{T_i}, A_{T_j}] \neq 0$), reproducing non-Abelian structure (Section~\ref{sec:edge_split}).
\end{enumerate}

An exploratory gravitational model (Surface Tension Gravity) is presented in Appendix~A. While STG reproduces GR light deflection and orbital precession at 1PN order---with the metric exponent $\alpha = 1/2$ derived from 600-cell topology---its strong-field predictions (notably a black hole shadow at half the GR size) are in $\sim 7\sigma$ tension with EHT observations. The gravitational sector should therefore be regarded as a toy model requiring substantial refinement, rather than a mature prediction of the framework.

\textbf{Limitations and open problems:}
\begin{itemize}
    \item The DSI mechanism (Section~\ref{sec:dsi}) explains \emph{why} masses scale as $\phi^n$, and complexity minimization (Section~\ref{sec:ab_derivation}) determines $(a,b)$ uniquely. The level selection (Section~\ref{sec:level_selection}) achieves 31/31 accuracy with per-line refinements: norm positivity ($N > 0$) on the $b=1$ line is derived, but the three-generation limit ($b \leq 2$ on $a=1$) remains a pattern. A systematic investigation (EXP-127) of 10 approaches---including norm analysis, $H_4$ representation theory, $E_8$ branching, and diameter constraints---found no geometric derivation of the generation count.
    \item The CKM exponents $\{3, 7, 12\}$ are expressible through $H_4$ Coxeter exponents (Section~6.2) but the connection is not uniquely determined. PMNS exponents ($n = 0, 1, 4$) are motivated by geometric dimensions.
    \item The gauge group derivation is now complete (EXP-135): $D_4$ from 24-cell, triality broken by $E_8$ Dynkin topology, leg assignment via Kac labels $\{3,4,5\}$, and $SU(2)_R \to U(1)_Y$ from dead-end isolation. The remaining open questions are the breaking \emph{scale} (the $W_R$ mass) and hypercharge normalization.
    \item \textbf{Neutrino masses} are not addressed by the current framework. The second 600-cell $T = \phi' \cdot S$ in the $E_8$ construction does not produce viable neutrino masses: the Galois identity $\phi \cdot \phi' = -1$ kills level-dependent mass formulas, and Type-I seesaw with $m_D = m_{\text{lepton}}$ gives a hierarchy too steep by a factor $\sim 2450$ (EXP-134). The neutrino sector likely requires a different mechanism (Majorana masses, radiative generation, or the $(1,8)$ sector of $E_8$ branching).
    \item The Down quark mass prediction (5.15\,MeV) lies at the 1$\sigma$ upper bound of the PDG value ($4.67^{+0.48}_{-0.17}$\,MeV), placing it within $1\sigma$ of experiment (EXP-128). The Down--Strange tension is structural and irreducible within shared corrections.
    \item The $E_8$ unification uses MSSM beta functions as a proxy; a direct derivation from $E_8 \to H_4$ breaking is needed.
    \item The link between discrete graph properties (edge counting) and continuous metric components (Appendix~A) remains heuristic; a rigorous derivation from the Lagrangian is needed.
    \item \textbf{Coupling constant scales (EXP-150):} The ratio $\alpha_s/\alpha = 10\phi$ uses mixed renormalization scales ($\alpha$ at $q = 0$ Thomson limit, $\alpha_s$ at $M_Z$). At a single scale $M_Z$, the ratio is 15.09 (6.8\% off from $10\phi$). The framework encodes coupling constants at their natural measurement scales, consistent with a non-perturbative geometric origin rather than perturbative running.
    \item \textbf{Electroweak scale (EXP-152):} The framework derives mass \emph{ratios} ($m_H/m_W$, $m_Z/m_W$) but not the absolute electroweak scale $v = 246\,$GeV, which remains an additional input. The top Yukawa coupling $y_t = 0.992 \approx 1$ and $n_{\text{top}} = 26 = \dim(\text{phi-sector})$ suggest a special role for the top quark, but no derivation of $v$ from the 600-cell has been found.
\end{itemize}

Despite these limitations, the framework offers a concrete, falsifiable geometric hypothesis for the origin of fundamental constants. The precision of the coupling constant derivations ($\alpha$ at $0.0001\%$, $\alpha_s$ at $0.03\%$) and the lepton mass hierarchy ($0.25\%$ after corrections) suggest that the 600-cell geometry captures genuine structure of physical law, even if the complete picture remains incomplete.

%==============================================================
\section*{Acknowledgments}
%==============================================================

We acknowledge the foundational work on spectral geometry by Alain Connes and Ali Chamseddine \cite{chamseddine1997}, the geometric optimality results of Cohn and Kumar \cite{cohn2007,cohn2011}, and the 600-cell GUT framework and Surface Tension Gravity developed by Christopher C. O'Neill \cite{oneill2025gut,oneill2025stg,oneill2025120cell,oneill2023binary}. We also acknowledge the experimental data from CODATA 2022 \cite{codata2022}, the Particle Data Group \cite{pdg2024}, the Fermi LAT collaboration \cite{fermi2009}, and the Event Horizon Telescope collaboration \cite{eht2019}.

%==============================================================
\appendix
%==============================================================

\section{Surface Tension Gravity: An Exploratory Model}
\label{app:stg}

\textbf{Note:} This appendix presents an exploratory gravitational model derived from the 600-cell compression factor $C(r)$. While the weak-field results (light deflection, orbital precession) match GR at first post-Newtonian order, the strong-field predictions---particularly the black hole shadow---are in significant tension with observations. This model should be regarded as a toy model that illustrates how the 600-cell geometry \emph{could} constrain spacetime structure, rather than a definitive gravitational theory.

\subsection{The STG Metric}

The gravitational compression factor $C(r) = 1 + M/r$ (Section~4.2) defines a spacetime metric. The most general static, spherically symmetric form consistent with $C(r)$ involves a single parameter $\alpha$:

\begin{equation}
ds^2 = -\frac{dt^2}{C^2} + C^2\,dr^2 + C^{2\alpha}\,r^2\,d\Omega^2
\label{eq:stg_metric}
\end{equation}

where $C(r) = 1 + M/r$. Unlike Schwarzschild, this metric has \textbf{no event horizon}: $g_{tt} = -1/C^2 < 0$ for all $r > 0$.

\subsection{Closed-Form Force Law}

From the geodesic equation in the Newtonian limit:

\begin{equation}
F(r) = \frac{Mr}{(r + M)^3}
\end{equation}

This recovers Newton ($F \approx M/r^2$) at $r \gg M$ but vanishes as $r \to 0$, preventing the classical singularity. The modified Kepler law $T^2 = 4\pi^2(r + M)^3/M$ has been verified numerically ($n = 3.0002$, error $0.007\%$).

\subsection{Weak-Field Classical Tests}

\textbf{Light deflection:} $\Delta\phi_{\text{STG}} = 4M/b$ (identical to GR), confirmed by numerical ODE integration.

\textbf{Gravitational redshift:} $z = M/r$ (matches GR at weak field).

\textbf{Orbital precession:} The perihelion precession satisfies $\Delta\phi_{\text{STG}}/\Delta\phi_{\text{GR}} = 1/2 + \alpha$ at 1PN order (Eq.~\ref{eq:precession_ratio}). For $\alpha = 1/2$, STG exactly reproduces the GR result $\Delta\phi = 6\pi M/r_0$.

\begin{equation}
\frac{\Delta\phi_{\text{STG}}}{\Delta\phi_{\text{GR}}} = \frac{1}{2} + \alpha + \mathcal{O}(M/r_0)
\label{eq:precession_ratio}
\end{equation}

\subsection{Geometric Origin of $\alpha = 1/2$}

For a given vertex $v_0$ on the 600-cell graph, BFS assigns each vertex a graph distance $d$. An edge is \emph{angular} if both endpoints share the same $d$, and \emph{radial} if $|d(v_i) - d(v_j)| = 1$.

\textbf{Result:} Exactly 360 of the 720 edges are angular for \emph{any} choice of $v_0$ ($\sigma = 0$ across all 120 positions):

\begin{equation}
\boxed{\frac{N_{\text{angular}}}{N_{\text{total}}} = \frac{360}{720} = \frac{1}{2}}
\label{eq:angular_half}
\end{equation}

This topological invariant motivates $\alpha = 1/2$. However, the per-shell angular fraction varies significantly:

\begin{table}[H]
\centering
\begin{tabular}{cccl}
\toprule
\textbf{Shell $d$} & \textbf{Angular frac.} & \textbf{Status} & \textbf{Note} \\
\midrule
1 & 0.294 & Invariant ($\sigma = 0$) & Near mass: mostly radial \\
2 & 0.405 & Invariant ($\sigma = 0$) & \\
3 & 0.476 & Invariant ($\sigma = 0$) & \\
4 & 0.909 & Invariant ($\sigma = 0$) & Far from mass: mostly angular \\
\bottomrule
\end{tabular}
\caption{Per-shell angular fraction from EXP-106. All values are topological invariants.}
\end{table}

\textbf{Caveat:} The link between edge-counting on the graph and the continuous metric exponent is heuristic. A rigorous derivation from the spectral action or an equation of motion is needed.

\subsection{Strong-Field Predictions and EHT Tension}

\begin{table}[H]
\centering
\begin{tabular}{lcc}
\toprule
\textbf{Observable} & \textbf{STG ($\alpha = 1/2$)} & \textbf{GR (Schwarzschild)} \\
\midrule
Event horizon & None ($g_{tt} < 0\;\forall\;r>0$) & $r = 2M$ \\
ISCO & $r = \tfrac{1+\sqrt{3}}{2}\,M \approx 1.37\,M$ & $r = 6M$ \\
Photon sphere & $r = M/2$ & $r = 3M$ \\
Shadow ($b_c$) & $\tfrac{3\sqrt{3}}{2}\,M \approx 2.60\,M$ & $3\sqrt{3}\,M \approx 5.20\,M$ \\
Weak-field deflection & $4M/b$ & $4M/b$ \\
Weak-field precession & $6\pi M/r_0$ & $6\pi M/r_0$ \\
\bottomrule
\end{tabular}
\caption{Comparison of STG and GR predictions.}
\label{tab:stg_vs_gr}
\end{table}

The shadow ratio $b_{c,\text{STG}}/b_{c,\text{GR}} = 1/2$ is exact. The EHT measured the M87$^*$ shadow at $42 \pm 3\,\mu\text{as}$ \cite{eht2019}, consistent with GR. STG predicts $\sim 21\,\mu\text{as}$, a $\sim 7\sigma$ discrepancy.

An analysis of the 600-cell lattice Green's function (EXP-106) shows that higher-order corrections to $C(r)$ from the discrete structure have the \emph{wrong sign} to resolve this tension. The lattice Green's function gives $C(r) \approx 1 + a_1 M/r + a_2 (M/r)^2$ with $a_2/a_1 < 0$, which makes the shadow \emph{smaller}, not larger. This constitutes a genuine strong-field failure of the STG model in its current form.

%==============================================================
\begin{thebibliography}{99}

% === 600-Cell Theory & Surface Tension Gravity (STG) ===

\bibitem{oneill2025gut}
C. C. O'Neill, ``The 600-Cell Grand Unified Theory: A Geometric and Group-Theoretic Model of the Standard Model and Beyond,'' preprint (2025).

\bibitem{oneill2025120cell}
C. C. O'Neill, ``The 120-Cell Theory of Everything,'' preprint (2025).

\bibitem{oneill2025stg}
C. C. O'Neill, ``Geometry-Primary Extension of Surface-Tension Gravity: An Effective Spacetime Metric from Compression $C(r)$ and the Classical Solar-System Tests,'' preprint (2025).

\bibitem{oneill2025scaffold}
C. C. O'Neill, ``A Rotation-Bundle Scaffold for EM-STG with SU(3) and an Electroweak Interface (with an STG-Higgs portal),'' preprint (2025).

\bibitem{oneill2025gps}
C. C. O'Neill, ``Surface-Tension Gravity Calibrated by GPS Time Dilation,'' preprint (2025).

\bibitem{oneill2023binary}
C. C. O'Neill, ``Isomorphism of a Modified Standard Model Particle Assignment with the Binary Octahedral Group (24-cell),'' preprint (2023).

% === Geometric Optimality & Spectral Principles ===

\bibitem{cohn2007}
H. Cohn and A. Kumar, ``Universally optimal distribution of points on spheres,'' \textit{J. Amer. Math. Soc.} \textbf{20}, 99--148 (2007).

\bibitem{chamseddine1997}
A. H. Chamseddine and A. Connes, ``The Spectral Action Principle,'' \textit{Commun. Math. Phys.} \textbf{186}, 731--750 (1997).

\bibitem{cohn2011}
H. Cohn and J. Woo, ``Three-point bounds for energy minimization,'' \textit{J. Amer. Math. Soc.} (2011).

% === Experimental Data ===

\bibitem{codata2022}
CODATA, ``Recommended Values of the Fundamental Physical Constants,'' NIST (2022).

\bibitem{fermi2009}
Fermi LAT Collaboration, ``A limit on the variation of the speed of light arising from quantum gravity effects,'' \textit{Nature} \textbf{462}, 331--334 (2009).

\bibitem{pdg2024}
Particle Data Group, ``Review of Particle Physics,'' \textit{Phys. Rev. D} \textbf{110}, 030001 (2024).

\bibitem{eht2019}
Event Horizon Telescope Collaboration, ``First M87 Event Horizon Telescope Results. I. The Shadow of the Supermassive Black Hole,'' \textit{Astrophys. J. Lett.} \textbf{875}, L1 (2019).

% === Lattice & Number Theory ===

\bibitem{conway1998icosian}
J. H. Conway and N. J. A. Sloane, ``Sphere Packings, Lattices and Groups,'' 3rd ed., Springer (1998). Ch.~8: The icosian lattice and $E_8$.

% === Discrete Scale Invariance ===

\bibitem{sornette1998}
D. Sornette, ``Discrete-scale invariance and complex dimensions,'' \textit{Phys. Rep.} \textbf{297}, 239--270 (1998).

\end{thebibliography}

\end{document}
